% Appendix ready to be included in the thesis.
% Required packages in the thesis preamble:
% booktabs, longtable, tabularx, array, enumitem, amsmath, amssymb,
% graphicx, xurl, hyperref, natbib (or an equivalent citation package).
%
% % This protocol is written for the complete study interval. Operational batching
% does not alter the methodological scope, questions, eligibility rules, or synthesis.

\providecommand{\protocolfield}[1]{\textbf{[#1]}}
\providecommand{\yes}{\textit{Yes}}
\providecommand{\partial}{\textit{Partially}}
\providecommand{\no}{\textit{No}}
\newcommand{\reviewtitle}{Quality constructs, deviations, and measures in LLM-generated UML diagrams}

\chapter{Protocol for the systematic mapping study on \reviewtitle}
\label{app:quality_focused_mapping_protocol}

\section{Introduction}
\label{sec:protocol_identification}

This appendix presents the protocol for the systematic mapping study that structures and synthesizes quality evidence in LLM-based UML diagram generation. The study aims to identify primary studies on the generation, transformation, completion, repair, refinement, or revision of UML diagram content by LLMs from textual software requirements or explicit textual domain descriptions. It further aims to analyze the extractable evidence those studies report about syntactic and semantic quality, syntax-semantic dissonance, evaluation references, measurement procedures, pragmatic adequacy, and correction effort.

The protocol presents the planning, search, selection, extraction, and synthesis principles established from the practice of secondary studies in Software Engineering \cite{Kitchenham2007Guidelines,Brereton2007Lessons,Petersen2015Mapping,Zhang2011Identifying}.

\begin{table}[htbp]
\centering
\small
\caption{Protocol identification}
\label{tab:protocol_identification}
\begin{tabularx}{\textwidth}{p{0.28\textwidth} X}
\toprule
\textbf{Field} & \textbf{Specification} \\
\midrule
Review title & \reviewtitle: a systematic mapping study \\
Review type & Focused systematic mapping study with structured extraction and synthesis of quality evidence \\
Protocol version & 3.0 \\
Protocol date & 2026-08-21 \\
Review team, supervision, and adjudication & Stated once, with roles and responsibilities, in Table~\ref{tab:protocol_reviewers} \\
Registration repository & \protocolfield{OSF, Zenodo, or institutional repository identifier} \\
Planned search interval & January 2022 to the final search date \\
Language of eligible reports & English \\
Primary disciplinary domains & Software Engineering, Requirements Engineering, Software Modeling, Model Driven Engineering, and Large Language Models \\
Protocol status & Prospective protocol. Selection, extraction, and synthesis must not begin before protocol validation and registration, except for explicitly designated pilot activities. \\
\bottomrule
\end{tabularx}
\end{table}

\section{Rationale and need for the secondary study}
\label{sec:protocol_rationale}

Recent primary studies indicate that LLMs can generate UML diagrams that satisfy at least some criteria of textual validity or syntactic conformity while presenting problems of semantic correctness or completeness relative to source requirements and domain information \cite{DeBari2024UML,Ferrari2024Sequence,Garaccione2025Class,Garaccione2025UseCase,Xiao2025Sequence}. These studies use different diagram types, source specifications, model families, prompting strategies, evaluation references, units of analysis, metrics, and terminology. Consequently, apparently similar labels such as syntactic correctness, semantic correctness, completeness, consistency, hallucination, or understandability may denote non equivalent constructs and procedures.

The conceptual background on model quality distinguishes the relation between a model and its language, the relation between a model and the relevant domain, and the relation between a model and its interpreters or uses \cite{Lindland1994Quality,Krogstie2006Process}. The UML specification also distinguishes concrete representations from the abstract syntax and well-formedness constraints of the modeling language \cite{OMG2017UML}. These distinctions are necessary for evaluating generated diagrams, but the manner in which recent LLM studies operationalize them remains to be systematically reconstructed.

Related secondary studies provide important but different forms of coverage. Reviews of UML model quality predate the current LLM-based generation setting \cite{Genero2011UMLQuality}. Reviews of transformations from natural language to UML mainly cover rule-based and earlier machine learning approaches \cite{Ahmed2022NaturalUML}. Recent studies map LLM applications across Model Driven Engineering or review UML and entity relationship diagram generation at a broad level \cite{Zhang2027LLMMDE,Rahmanian2026Diagrams}. None of them maps LLM-based UML generation studies that report extractable quality evidence, and none synthesizes how such studies define, classify, measure, and validate diagram quality and the inadequacies they report.

To establish whether such a secondary study already exists, a verification search was run on 2026-08-20 against the OpenAlex REST API, over the titles, abstracts, and keywords of works published from January 2022. OpenAlex was used for this check because it indexes the catalogues of the publishers relevant to this review, including IEEE, ACM, Springer, and Elsevier, and because it is queryable without a subscription or an institutional session, so that a reader can repeat the check from the string and the date alone. Its purpose was not to build the corpus of this review but to answer a single prior question: whether a review, mapping study, or survey of LLM-based UML generation already reports what this protocol sets out to report. The string is given in Table~\ref{tab:protocol_verification_search}. It pairs the population and intervention terms of the main search with the publication-type terms by which a secondary study identifies itself, so that an existing review on this topic would be returned even if it used a different vocabulary for quality.

The bare phrase \textit{model driven} is deliberately absent from the modelling block, which uses \textit{model-driven engineering} instead. As a bare phrase it matches the compound \textit{large language model-driven}, a productive construction in fields unrelated to software modelling; retaining it returned 308 records rather than 111, and the surplus consisted almost entirely of works in which a language model drives something that is not a model.

\begin{table}[htbp]
\centering
\caption{Verification search string, as submitted to the OpenAlex REST API on 2026-08-20}
\label{tab:protocol_verification_search}
\begin{tabular}{@{}l@{}}
\toprule
\begin{minipage}{0.95\textwidth}\vspace{3pt}
\scriptsize\ttfamily\raggedright
filter=title\_and\_abstract.search:\\
\quad ( ``large language model'' OR ``large language models'' OR LLM OR LLMs OR\\
\quad ``generative AI'' OR ChatGPT OR GPT )\\
\quad AND\\
\quad ( UML OR ``unified modeling language'' OR PlantUML OR\\
\quad ``class diagram'' OR ``use case diagram'' OR ``sequence diagram'' OR\\
\quad ``activity diagram'' OR ``state machine diagram'' OR\\
\quad ``model-driven engineering'' )\\
\quad AND\\
\quad ( ``systematic review'' OR ``systematic literature review'' OR\\
\quad ``systematic mapping'' OR ``mapping study'' OR ``literature review'' OR\\
\quad ``secondary study'' OR ``tertiary study'' OR survey OR\\
\quad ``state of the art'' )\\
, from\_publication\_date:2022-01-01
\vspace{3pt}
\end{minipage}\\
\bottomrule
\end{tabular}
\end{table}

The search returned 111 records, which reduce to 79 distinct works once preprint, published, and repository versions of one title are grouped. Screening titles and abstracts identified four secondary studies on LLM-based generation of software models. Two are those discussed above \cite{Zhang2027LLMMDE,Rahmanian2026Diagrams}. The other two are a survey of large language models for Model-Driven Engineering and an anonymized systematic mapping study, still under review, of large language models for the generation, extraction, and validation of UML diagrams.

The survey covers Model-Driven Engineering as a whole rather than the quality of generated diagrams. The mapping study covers generation together with extraction from code, and its public extraction sheet (\texttt{10.5281/zenodo.21343580}) records each study's evaluation strategy but not the construct measured, the reference it was measured against, or the inadequacies reported; its evaluation type labels are distinct for every one of its twenty-four extracted studies. Neither reconstructs the quality constructs, which is the operation this review adds. One record matched without being a secondary study, a student survey on LLMs and UML modelling, returned because \textit{survey} belongs to the publication-type block. The returned records, the string, the request URL, and the date are kept with the review materials.

The review is necessary because the subsequent stages of the thesis depend on that account. Before the incidence of inadequacies can be estimated, their categories and units must be defined. Before contextual determinants can be investigated, the dependent quality constructs must be operationalized. Before pragmatic effects or rework can be evaluated, the literature must be examined for the tasks, measures, and procedures that already exist. The secondary study therefore has both a descriptive and a constitutive function: it maps the field, and it establishes the evidence base required to formulate the later empirical questions and instruments.

\section{Researchers involved}
\label{sec:protocol_reviewers}

The search, study selection, and synthesis stages are carried out by researcher Helaine Barreiros. Study selection, data extraction, and the application of the data collection instrument are conducted with the participation of researcher Emanoel Barreiros. The process is supervised by Professor Ivaldir Hon\'orio de Farias J\'unior and co-supervised by Professor Cleyton M\'ario de Oliveira Rodrigues, who jointly resolve disagreements and validate the method. Domain expertise is distributed as follows: Ivaldir Hon\'orio de Farias J\'unior in Software Engineering, Cleyton M\'ario de Oliveira Rodrigues in Artificial Intelligence, Adauto Trigueiro de Almeida Filho in Software Engineering, UML, and empirical Software Engineering, and Helaine Barreiros and Emanoel Barreiros in empirical Software Engineering, UML, and Software Engineering. Table~\ref{tab:protocol_reviewers} states the role of each researcher.

\begin{table}[htbp]
\centering
\small
\caption{Reviewer roles}
\label{tab:protocol_reviewers}
\begin{tabularx}{\textwidth}{p{0.22\textwidth} p{0.26\textwidth} X}
\toprule
\textbf{Role} & \textbf{Responsible researcher} & \textbf{Responsibilities} \\
\midrule
Primary reviewer & Helaine Barreiros & Search execution, record management, screening, extraction, synthesis \\
Second reviewer & Emanoel Barreiros & Independent calibration, sampled title and abstract screening, complete full-text screening, quality-evidence extraction, protocol validation \\
Adjudicator & Ivaldir Hon\'orio de Farias J\'unior and Cleyton M\'ario de Oliveira Rodrigues & Resolution of unresolved eligibility, extraction, appraisal, and coding disagreements \\
Methodological validator & Ivaldir Hon\'orio de Farias J\'unior and Cleyton M\'ario de Oliveira Rodrigues & Review of protocol, search strategy, selection process, synthesis, and reporting alignment \\
UML domain validator & Adauto Trigueiro de Almeida Filho & Review of diagram terminology, codebook boundaries, UML carrier categories, and reference distinctions \\
LLM domain validator & Cleyton M\'ario de Oliveira Rodrigues & Review of LLM terminology, model configuration fields, and context categories \\
\bottomrule
\end{tabularx}
\end{table}


Every decision remains attributable to an individual reviewer and date. Consensus decisions do not overwrite original independent decisions. The audit trail preserves both. Reviewers declare prior authorship or close collaboration with any candidate primary report. A reviewer does not make the final eligibility or appraisal decision for their own work.

\section{Study scope}
\label{sec:protocol_scope}

The study scope comprises scientific reports that investigate LLM-based generation of UML diagrams from textual software requirements or explicit textual domain descriptions. Three units are distinguished throughout the review. Table~\ref{tab:protocol_units} states what each one is and the role it plays.

\begin{longtable}{p{0.20\textwidth} p{0.40\textwidth} p{0.32\textwidth}}
\caption{Units of selection, counting, and extraction}
\label{tab:protocol_units}\\
\toprule
\textbf{Unit} & \textbf{What it is} & \textbf{Role in the review} \\
\midrule
\endfirsthead
\toprule
\textbf{Unit} & \textbf{What it is} & \textbf{Role in the review} \\
\midrule
\endhead
Report & The document screened for eligibility. & Unit of selection. \\
Primary study & The empirical or technical investigation, which may be reported in one or more publications. & Unit used to prevent double counting. \\
Evaluation instance & An extractable combination of diagram type, input corpus, model configuration, generation condition, and evaluation procedure. & Preferred unit for analytical extraction when a report contains more than one of these. \\
\bottomrule
\end{longtable}

The study focuses on primary studies published in English between January 2022 and July 2026. Journal articles, conference papers, and workshop papers are eligible when they provide a complete scientific report. Preprints are eligible because of the rapid evolution of the research area, but publication status is coded explicitly and peer reviewed and preprint evidence are analyzed separately when the distinction affects interpretation. Secondary and tertiary studies are excluded from the primary corpus but retained for background, search validation, and snowballing.

\section{Operational definitions}
\label{sec:protocol_definitions}

\begin{longtable}{p{0.22\textwidth} p{0.70\textwidth}}
\caption{Operational definitions used in the protocol}
\label{tab:protocol_definitions}\\
\toprule
\textbf{Term} & \textbf{Operational definition} \\
\midrule
\endfirsthead
\toprule
\textbf{Term} & \textbf{Operational definition} \\
\midrule
\endhead
Large Language Model & A neural language model identified by the study as a large, foundation, generative, instruction following, or conversational language model and used as a substantial component of the generation process. Rule based NLP and conventional machine learning without an LLM are outside scope. The operational boundary in Section~\ref{sec:substantive_llm_use_boundary} governs cases in which a report uses broad labels such as language model, pretrained language model, or foundation model without clearly establishing a generative or semantically substantive role in UML content production. \\
UML diagram & A model or diagram explicitly identified as UML or as one of the UML diagram types. PlantUML output is eligible only when it is intended to encode UML. Generic Mermaid, C4, ER, BPMN, or architectural diagrams are outside scope unless the UML component is separable. \\
Textual software requirements & Natural language or structured textual descriptions of intended software behavior, structure, goals, actors, rules, constraints, scenarios, or domain information. Examples include requirements documents, use case narratives, user stories, problem statements, feature descriptions, and domain descriptions. \\
Generation task & Production, completion, transformation, repair, refinement, or iterative revision of a UML diagram in which the LLM materially determines the resulting model content. \\
Generation context & The complete information made available to the model for an evaluation instance, including instructions, requirements, domain information, UML knowledge, examples, retrieved material, tools, prior turns, and ordering. \\
Textual validity & Conformance of the concrete output string to the grammar accepted by the selected representation or rendering tool. It does not by itself establish UML conformity. \\
UML syntactic conformity & Conformance of the represented model to the applicable UML abstract syntax, admissible constructs, relations, and well-formedness constraints. \\
Semantic correctness & Degree to which represented propositions, elements, relations, constraints, states, events, or behaviors are supported by the relevant requirements and domain reference. \\
Semantic completeness & Degree to which relevant units in the defined semantic reference are adequately represented in the diagram. \\
Pragmatic adequacy & Degree to which the diagram supports a defined engineering activity for a specified actor and context. Comprehensibility may be one component but is not treated as a sufficient synonym. \\
Inadequacy & A reported or inferable discrepancy between the generated diagram and an explicitly identified language, semantic, or task reference. The term is stipulated here as the umbrella for a deviation on any of the three quality axes. It is neither the negation of \textit{pragmatic adequacy} above, which names the use axis alone, nor the \textit{adequacy} of natural language generation evaluation, which names the semantic axis alone. Section~\ref{sec:protocol_attribute_deviation_measure} gives the axis-specific names and states why \textit{defect}, \textit{fault}, and \textit{anomaly} are not adopted. The review preserves the source term and separately assigns a normalized code. \\
Semantic reference & The source against which correctness or completeness is assessed, such as requirements, explicit domain knowledge, an annotated reference model, a validated model, a gold standard, a rubric, or expert judgment. \\
Correction or rework & Human or automated activity required after initial generation to inspect, change, validate, or regenerate a diagram so that it satisfies a defined acceptance criterion. \\
Publication family & Two or more reports that describe the same or substantially overlapping primary study, dataset, experiment, or artifact. \\
\bottomrule
\end{longtable}

\subsection{Attributes, deviations, and measures}
\label{sec:protocol_attribute_deviation_measure}

The three synthesis questions that carry the constitutive function of this review operate at three different levels, and those levels are the ones of the measurement vocabulary of systems and software engineering rather than a distinction introduced here. An \textit{attribute} is a property of an entity that can be distinguished quantitatively or qualitatively. A \textit{measure} assigns a value to an attribute through a stated measurement method, and an \textit{indicator} is a measure that supports interpretation against decision criteria \cite{ISO15939,ISO25020,SWEBOK2024}. A \textit{deviation} is a departure of the entity from an explicitly identified reference. The review records, for every included study, a single spine.

\begin{equation}
\text{attribute} \rightarrow \text{deviation} \rightarrow \text{reference} \rightarrow \text{unit} \rightarrow \text{indicator} \rightarrow \text{procedure} \rightarrow \text{aggregation}.
\label{eq:protocol_quality_spine}
\end{equation}

SQ1 covers the first link, SQ2 the second, and SQ3 the remainder, restated as Equation~\ref{eq:protocol_metrology_chain} where the metric catalogue is specified. Keeping the links apart is the purpose of the decomposition. A primary study that reports having measured semantic correctness and gives a percentage has named an attribute and reported an aggregated value, and has usually left the deviation, the unit, and the denominator unstated. Recovering those is what this review adds, and it cannot be done by an instrument that stores the three levels in one field. Conflating a measurable concept with its measure is the specific failure the separation is designed to prevent.

The two normalization vocabularies are therefore duals, and they must not share labels. Every value of the normalized construct field names an attribute, and every value of the normalized discrepancy operation field names a departure from one. The correspondences are semantic completeness with omission; semantic correctness with unsupported addition, incorrect representation, and incorrect classification; consistency with inconsistency; and textual validity and UML syntactic conformity with the corresponding nonconformity, which the violated reference field records at the level at which it occurred. A deviation that decomposes into an operation and a carrier, such as an incorrect relation or an incorrect actor, is recorded as that pair and never as a construct label, because storing it in both vocabularies would report a deviation type as a quality dimension in the construct matrix of SQ1. Effort quantities such as correction effort and inspection effort are measures and are recorded in the rework fields, not as construct labels.

The duality is not an identity, and its asymmetries are findings rather than defects of the instrument. A study may report deviations without ever naming the attribute they depart from, in which case the attribute is recorded as not reported, and a study may name an attribute and report no deviation at all. SQ1 and SQ2 are reported over the same corpus so that both asymmetries remain visible.

The axis-specific name of a deviation follows the reference it violates, because quality here is a relation and not a property of the diagram alone \cite{Lindland1994Quality,Krogstie2006Process}. On axis L a deviation is a nonconformity to the grammar of the concrete representation, or to the abstract syntax and well-formedness constraints of the modeling language \cite{OMG2017UML}. On axis D it is an invalidity when the diagram asserts what the reference does not support, and an incompleteness when the reference contains what the diagram omits; these are the two semantic goals of the model quality frameworks this protocol adopts. On axis U it is an inadequacy in the strict sense, namely a failure to support the defined engineering activity for the specified actor and context. In process terms, axes L and D correspond respectively to verification against a specified reference and to validation against intended use \cite{ISO12207}. Syntax-semantic dissonance is accordingly not a further construct but the co-occurrence of conformity on axis L with deviation on axis D, and it is computed from the axes rather than asserted.

The umbrella term is stipulated rather than borrowed, and the alternatives are declined for stated reasons. \textit{Defect} and \textit{fault} presuppose a specification that the generated diagram is not required to have, commit to repair, and carry the fault to failure chain of systems in operation. \textit{Anomaly} is tied to a reporting and disposition process that this review does not run. \textit{Nonconformity} is exact on axis L and wrong on the other two \cite{ISO24765,IEEE1044}. The product quality models of the SQuaRE series characterize a software product in operation, and do not supply the constructs at issue here, which concern a specification artifact evaluated against a domain reference; they are cited for the measurement framework and not as the quality model of this review \cite{ISO25010,ISO25020}.


\section{Objective and expected contribution}
\label{sec:protocol_objective}

The study objective is to map primary studies that use LLMs to generate, transform, complete, repair, refine, or revise UML diagram content from textual software requirements or explicit textual domain descriptions, to synthesize how these studies characterize, classify, operationalize, measure, validate, and report the syntactic, semantic, and pragmatic quality of the generated UML content, the inadequacies they report, and the evidence of syntax-semantic dissonance, and to establish the quality constructs, deviation categories, units of assessment, evaluation references, reliability practices, baselines, and correction effort measures required to formulate the empirical questions and instruments of the subsequent thesis stages.

The study is expected to produce six integrated contributions.

\begin{enumerate}[label=\roman*)]
    \item a structured map of publication, study, diagram, task, input, model, context, output, and evaluation characteristics;
    \item a traceable matrix connecting quality terms to definitions, theoretical references, evaluation references, units, indicators, metrics, and procedures;
    \item a normalized taxonomy of inadequacies reported in the literature, while preserving the original terminology used by each primary study;
    \item a catalogue of semantic references, model oracles, human evaluation procedures, automated checks, and reliability practices;
    \item a synthesis of reported contextual conditions, pragmatic effects, and correction or rework measures;
    \item a specification of the quality constructs, deviation categories, units of assessment, evaluation references, baselines, and correction effort measures for which the literature provides precedent, together with an explicit record of those for which it does not, as the basis for the empirical instruments of the subsequent thesis stages.
\end{enumerate}

\section{Research questions}
\label{sec:protocol_questions}

\begin{quote}
\textbf{ORQ.} How have primary studies used LLMs to generate, transform, complete, repair, refine, or revise UML diagram content from textual software requirements or explicit textual domain descriptions, under which generation context, and how do they characterize, classify, operationalize, measure, validate, and report the syntactic, semantic, and pragmatic quality of the generated UML content, the reported inadequacies, and the evidence of syntax-semantic dissonance?
\end{quote}

The question is decomposed into mapping questions and synthesis questions.

\subsection{Mapping questions}

\begin{description}[style=nextline,leftmargin=2.0cm,labelwidth=1.5cm]
    \item[MQ1.] How is the research landscape distributed by publication year, venue, publication type, research method, contribution type, author affiliation, and availability of replication material?
    \item[MQ2.] Which UML diagram types, generation, transformation, completion, repair, refinement, or revision tasks, textual input forms, application domains, datasets, and output representations, including PlantUML, are investigated?
    \item[MQ3.] Which LLM families, model versions, access modes, prompting and context strategies, demonstrations, retrieval mechanisms, fine tuning procedures, agent structures, tool integrations, feedback mechanisms, and inference parameters are reported?
    \item[MQ4.] Against which baseline conditions do the mapped studies evaluate the generated UML content, how many baselines does a study use, and how many studies declare no baseline as opposed to not reporting one?
    \item[MQ5.] How do diagram types, generation configurations, and evaluation practices co occur across the literature, and which combinations remain under investigated?
\end{description}

\subsection{Quality evidence and synthesis questions}

\begin{description}[style=nextline,leftmargin=2.0cm,labelwidth=1.5cm]
    \item[SQ1.] Which quality constructs are used, how are they named and defined by the primary studies themselves, and how can the reported terms be normalized across the syntactic, semantic, and pragmatic dimensions without erasing study specific distinctions?
    \item[SQ2.] Which inadequacies are reported, how are they defined, against which violated reference is each one a departure, which UML elements or relations carry them, and how can the original categories be normalized without erasing diagram specific distinctions?
    \item[SQ3.] Which metrics, rubrics, oracles, evaluation references, units of assessment, scales, thresholds, aggregation rules, automated tools, and judgment procedures are used?
    \item[SQ4.] To what extent does the reported evidence distinguish syntactic quality from semantic quality, and what direct or indirect evidence of syntax-semantic dissonance is available?
    \item[SQ5.] How do the studies establish the reliability and validity of quality assessments, including evaluator training, independent coding, agreement measurement, adjudication, repeated generation, sensitivity analysis, and reporting of uncertainty?
    \item[SQ6.] How is the knowledge required for generation supplied and structured as generation context, with respect to the domain of the specified system and to the technical contract of the UML notation?
    \item[SQ7.] How are pragmatic adequacy, inspection, correction, rework, or productivity effects defined and measured, and with what inferential strength are they claimed?
\end{description}

\subsection{Traceability between questions, data, and synthesis}

\begin{longtable}{@{}p{0.06\textwidth} p{0.25\textwidth} p{0.24\textwidth} p{0.25\textwidth}@{}}
\caption{Traceability from research questions to extracted data and synthesis products}
\label{tab:protocol_rq_traceability}\\
\toprule
\textbf{RQ} & \textbf{Required data} & \textbf{Synthesis procedure} & \textbf{Primary product} \\
\midrule
\endfirsthead
\toprule
\textbf{RQ} & \textbf{Required data} & \textbf{Synthesis procedure} & \textbf{Primary product} \\
\midrule
\endhead
MQ1 & Bibliographic metadata, study design, contribution type, affiliation, artifact availability & Frequencies, temporal trends, venue classification, and cross tabulation & Publication and evidence landscape \\
MQ2 & Diagram type, task, input, domain, dataset, output representation & Classification and co-occurrence analysis & Diagram, task, and data map \\
MQ3 & Model, version, access, context, prompt, retrieval, adaptation, tools, parameters & Frequency distributions and configuration profiles & Technical configuration map \\
MQ4 & Baseline condition, one occurrence per baseline & Frequency distribution and comparison profile & Baseline and comparison map \\
MQ5 & Combined categories from MQ2 to MQ4 & Co occurrence matrices and gap analysis & Research concentration and gap map \\
SQ1 & Original terms, definitions, citations, operational definitions, examples & Conceptual comparison and normalized construct matrix & Quality construct matrix \\
SQ2 & Original inadequacy labels, textual definitions, examples, affected UML carrier, reference violated & Hybrid deductive and inductive coding, constant comparison, and negative case analysis & Taxonomy of reported inadequacies \\
SQ3 & Reference, oracle, unit, formula, denominator, scale, threshold, aggregation, procedure & Metrological comparison and traceability analysis & Evaluation reference and metric catalogue \\
SQ4 & Syntactic and semantic evidence availability, dissonance classification, evidence basis & Extraction-based comparison and evidence grouping & Syntax-semantic dissonance evidence map \\
SQ5 & Reviewer roles, training, independence, agreement, adjudication, repetitions, uncertainty & Methodological appraisal and reliability profile & Assessment credibility matrix \\
SQ6 & Domain knowledge supplied and UML technical instruction supplied & Classification and cross tabulation over the whole corpus & Generation context knowledge map \\
SQ7 & Context factor, inferential status, task, actor, outcome, correction and rework measures & Narrative synthesis stratified by inferential strength & Pragmatic and rework evidence map \\
\bottomrule
\end{longtable}

\section{Search process}
\label{sec:protocol_sources}

\subsection{Information sources}
\label{sec:protocol_information_sources}

The identification strategy combines broad bibliographic indexes, direct publisher or professional society libraries, citation searching, and manual venue inspection. These route types perform complementary functions and are not interchangeable: broad indexes reduce dependence on a predefined venue list, direct libraries improve access to conference and workshop proceedings, citation searching reaches studies whose terminology the query does not cover, and manual venue inspection verifies coverage gaps rather than defining the documentary population. The configuration follows methodological guidance on complementary search routes in Software Engineering secondary studies \cite{Brereton2007Lessons,Zhang2011Identifying,Wohlin2014Snowballing} and recent secondary studies showing that LLM-based modeling research is distributed across main conferences, companion and workshop tracks, and journals \cite{Zhang2027LLMMDE,Rahmanian2026Diagrams}.

\begin{table}[htbp]
\centering
\scriptsize
\caption{Information sources, status, and function}
\label{tab:protocol_sources}
\begin{tabularx}{\textwidth}{p{0.20\textwidth} p{0.15\textwidth} X}
\toprule
\textbf{Source} & \textbf{Status} & \textbf{Function and treatment} \\
\midrule
Scopus & Executed & Broad multidisciplinary index. Title, abstract, and keyword search with complete export and preserved provenance. \\
IEEE Xplore & Executed & Direct coverage of IEEE and IEEE/ACM venues. Metadata search adapted to platform syntax. \\
ACM Digital Library & Executed & Direct coverage of ACM venues, over the ACM Full-Text Collection. \\
Citation searching & Executed as a complementary route & Backward and forward citation searching seeded from the included set and from the seed studies of Table~\ref{tab:protocol_seed_set}. \\
Manual venue inspection & Executed as coverage validation & Tables of contents of the venues of Table~\ref{tab:protocol_manual_venues}. Does not replace the automated search. \\
Web of Science; ScienceDirect; arXiv & \protocolfield{Decision pending} & Candidate complementary sources inherited from the earlier protocol version and not executed for the present corpus. Each is either activated or its exclusion is justified before the search is declared closed. \\
\bottomrule
\end{tabularx}
\end{table}

The automated search over the three executed indexes defines the candidate documentary corpus. The other routes complement it, so the conclusions of the review are not restricted to a predetermined set of conferences or journals, and a study published outside the manual venue set is eligible whenever it satisfies the criteria of Section~\ref{sec:protocol_eligibility}.

\subsubsection{Prespecified seed set}

Five primary studies are prespecified as the seed set. The seed set is a validation instrument and not a privileged inclusion list: every seed study is subject to the same eligibility criteria as every other retrieved report, and a study identified after the search began is never retroactively added to it.

\begin{longtable}{p{0.06\textwidth} p{0.60\textwidth} p{0.26\textwidth}}
\caption{Prespecified seed studies used in search validation}
\label{tab:protocol_seed_set}\\
\toprule
\textbf{ID} & \textbf{Study} & \textbf{DOI} \\
\midrule
\endfirsthead
\toprule
\textbf{ID} & \textbf{Study} & \textbf{DOI} \\
\midrule
\endhead
SE1 & De Bari et al., \emph{Evaluating Large Language Models in Exercises of UML Class Diagram Modeling} \cite{DeBari2024UML} & \url{https://doi.org/10.1145/3674805.3690741} \\
SE2 & Ferrari et al., \emph{Model Generation with LLMs: From Requirements to UML Sequence Diagrams} \cite{Ferrari2024Sequence} & \url{https://doi.org/10.1109/REW61692.2024.00044} \\
SE3 & Garaccione et al., \emph{A Comparison of Different Large Language Models for the Generation of UML Class Diagrams} \cite{Garaccione2025Class} & \url{https://doi.org/10.1109/MODELS-C68889.2025.00078} \\
SE4 & Garaccione et al., \emph{Evaluating Large Language Models in Exercises of UML Use Case Diagrams Modeling} \cite{Garaccione2025UseCase} & \url{https://doi.org/10.1109/NLBSE66842.2025.00015} \\
SE5 & Xiao et al., \emph{UML Sequence Diagram Generation: A Multi-Model, Multi-Domain Evaluation} \cite{Xiao2025Sequence} & \url{https://doi.org/10.1109/ICSE-SEIP66354.2025.00030} \\
\bottomrule
\end{longtable}

Seed recovery is the acceptance test for a search string: a candidate string is accepted only when it retrieves the complete seed set, or when every failure has a documented complementary route. Because the five seeds use comparatively explicit UML terminology, seed recovery is treated as a necessary and not a sufficient condition, and it is never used on its own to justify narrowing a query block.

\subsubsection{Citation searching}

Citation searching is a route of its own and not an informal supplement, because the review deliberately searches without quality terminology and a study that names neither its diagrams nor its models in the query vocabulary is reachable only through the reference lists of studies that do. It is executed in two directions from a defined start set, namely the included studies and the seed studies of Table~\ref{tab:protocol_seed_set}. Backward searching inspects the reference list of every study in the start set. Forward searching inspects the works citing it, using the citation data available in the executed indexes.

Every record reached by citation searching enters the same record list and the same filters as a record from any other route, carries the route and the study that led to it as provenance, and is reported in its own PRISMA count. The route is iterated until one pass over a start set yields no new eligible record.

\subsubsection{Source contribution and retention analysis}

A source contribution matrix records, for every route, the number of records it identified, the number unique to it, and the number of those that survived screening. A route is retained when it contributes unique eligible records or resolves a documented coverage gap, and its removal or non-activation is justified in writing before the search is declared closed. Overlap between routes is not by itself a reason to remove one, because provenance-preserving deduplication already prevents double counting and the overlap is itself evidence about coverage. The analysis also records known misses. \textit{Automating Software Diagram Generation with Large Language Models} \cite{Gheorghita2025Diagrams}, published in an Elsevier venue and within the scope of this review, is retrieved by no executed source. It is the case that activation of the pending complementary sources would have to resolve, and it is recorded here so that the coverage gap of Table~\ref{tab:protocol_sources} is a named absence rather than a hypothesis.

\subsection{Automated database search}

The automated database search is the primary mechanism for defining the candidate documentary corpus, and the search strategy is intentionally broad enough to identify candidate studies for the focused review, and was organized into three conceptual blocks: (LLM terms) AND (UML and diagram terms) AND (generation and source terms).  

\subsubsection{Scopus search string}

\begin{quote}
\begin{minipage}{0.96\textwidth}
\scriptsize\ttfamily\raggedright
TITLE-ABS-KEY (\\
\quad ( ``large language model*'' OR LLM* OR ``generative AI'' OR\\
\quad ``generative artificial intelligence'' OR ChatGPT OR GPT OR Claude OR\\
\quad Gemini OR Llama OR Qwen OR DeepSeek )\\
)\\
AND\\
(\\
\quad TITLE-ABS (\\
\qquad UML OR ``Unified Modeling Language'' OR PlantUML OR\\
\qquad ``class diagram*'' OR ``use case diagram*'' OR ``sequence diagram*'' OR\\
\qquad ``activity diagram*'' OR ``state machine diagram*'' OR ``state diagram*'' OR\\
\qquad ``component diagram*'' OR ``deployment diagram*'' OR ``object diagram*'' OR\\
\qquad ``communication diagram*'' OR ``package diagram*'' OR\\
\qquad ``composite structure diagram*'' OR ``timing diagram*'' OR\\
\qquad ``interaction overview diagram*''\\
\quad )\\
\quad OR\\
\quad AUTHKEY (\\
\qquad UML OR ``Unified Modeling Language'' OR PlantUML OR\\
\qquad ``class diagram*'' OR ``use case diagram*'' OR ``sequence diagram*'' OR\\
\qquad ``activity diagram*'' OR ``state machine diagram*'' OR ``state diagram*'' OR\\
\qquad ``component diagram*'' OR ``deployment diagram*'' OR ``object diagram*'' OR\\
\qquad ``communication diagram*'' OR ``package diagram*'' OR\\
\qquad ``composite structure diagram*'' OR ``timing diagram*'' OR\\
\qquad ``interaction overview diagram*''\\
\quad )\\
)\\
AND\\
TITLE-ABS-KEY (\\
\quad ( generat* OR creat* OR construct* OR synthes* OR produc* OR transform* OR\\
\quad complet* OR repair* OR refine* OR ``natural language'' OR requirement* OR\\
\quad specification* OR ``user stor*'' OR scenario* OR ``problem statement*'' OR\\
\quad ``domain description*'' )\\
)
\end{minipage}
\end{quote}

\subsubsection{IEEE Xplore search string}

\begin{quote}
\begin{minipage}{0.96\textwidth}
\scriptsize\ttfamily\raggedright
(\\
\quad (\\
\qquad ``Document Title'':``large language model*'' OR\\
\qquad ``Document Title'':LLM OR ``Document Title'':LLMs OR\\
\qquad ``Document Title'':``generative AI'' OR\\
\qquad ``Document Title'':``generative artificial intelligence'' OR\\
\qquad ``Document Title'':ChatGPT OR ``Document Title'':GPT OR\\
\qquad ``Document Title'':Claude OR ``Document Title'':Gemini OR\\
\qquad ``Document Title'':Llama OR ``Document Title'':Qwen OR\\
\qquad ``Document Title'':DeepSeek OR\\
\qquad ``Abstract'':``large language model*'' OR ``Abstract'':LLM OR\\
\qquad ``Abstract'':LLMs OR ``Abstract'':``generative AI'' OR\\
\qquad ``Abstract'':``generative artificial intelligence'' OR\\
\qquad ``Abstract'':ChatGPT OR ``Abstract'':GPT OR ``Abstract'':Claude OR\\
\qquad ``Abstract'':Gemini OR ``Abstract'':Llama OR ``Abstract'':Qwen OR\\
\qquad ``Abstract'':DeepSeek OR\\
\qquad ``Author Keywords'':``large language model*'' OR\\
\qquad ``Author Keywords'':LLM OR ``Author Keywords'':LLMs OR\\
\qquad ``Author Keywords'':``generative AI'' OR\\
\qquad ``Author Keywords'':``generative artificial intelligence'' OR\\
\qquad ``Author Keywords'':ChatGPT OR ``Author Keywords'':GPT OR\\
\qquad ``Author Keywords'':Claude OR ``Author Keywords'':Gemini OR\\
\qquad ``Author Keywords'':Llama OR ``Author Keywords'':Qwen OR\\
\qquad ``Author Keywords'':DeepSeek\\
\quad )\\
\quad AND\\
\quad (\\
\qquad ``All Metadata'':UML OR ``All Metadata'':``Unified Modeling Language'' OR\\
\qquad ``All Metadata'':PlantUML OR ``All Metadata'':``class diagram'' OR\\
\qquad ``All Metadata'':``use case diagram'' OR ``All Metadata'':``sequence diagram'' OR\\
\qquad ``All Metadata'':``activity diagram'' OR\\
\qquad ``All Metadata'':``state machine diagram'' OR ``All Metadata'':``state diagram'' OR\\
\qquad ``All Metadata'':``component diagram'' OR ``All Metadata'':``deployment diagram'' OR\\
\qquad ``All Metadata'':``object diagram'' OR ``All Metadata'':``communication diagram'' OR\\
\qquad ``All Metadata'':``package diagram'' OR\\
\qquad ``All Metadata'':``composite structure diagram'' OR\\
\qquad ``All Metadata'':``timing diagram'' OR\\
\qquad ``All Metadata'':``interaction overview diagram''\\
\quad )\\
\quad AND\\
\quad (\\
\qquad ``All Metadata'':generate OR ``All Metadata'':generation OR\\
\qquad ``All Metadata'':create OR ``All Metadata'':creation OR\\
\qquad ``All Metadata'':construct OR ``All Metadata'':construction OR\\
\qquad ``All Metadata'':synthesis OR ``All Metadata'':produce OR\\
\qquad ``All Metadata'':transform OR ``All Metadata'':transformation OR\\
\qquad ``All Metadata'':complete OR ``All Metadata'':completion OR\\
\qquad ``All Metadata'':repair OR ``All Metadata'':refine OR\\
\qquad ``All Metadata'':refinement OR ``All Metadata'':``natural language'' OR\\
\qquad ``All Metadata'':requirement OR ``All Metadata'':specification OR\\
\qquad ``All Metadata'':``user story'' OR ``All Metadata'':scenario OR\\
\qquad ``All Metadata'':``problem statement'' OR\\
\qquad ``All Metadata'':``domain description''\\
\quad )\\
)
\end{minipage}
\end{quote}

\subsubsection{ACM Full-Text Collection search strategy v1.0}

The query must be implemented using three combined rows using \texttt{AND} in \emph{Search Within} interface. The first row is set to \emph{Anywhere} but contains explicit field prefixes so that the LLM block is evaluated across title, abstract, or author keywords. The UML block is deliberately not restricted in the same way. A validation round that restricted it to title, abstract, and author keywords was rejected, because differential inspection identified eligible model-generation studies that expose no UML terminology in those fields. The diagnostic false-negative case was \textit{Multi-step Iterative Automated Domain Modeling with Large Language Models} \cite{Yang2024MultiStepDomainModeling}, which the broad block retrieves and which the corpus retains. Restricting the block would have removed from the corpus a study of exactly the phenomenon under review, which is the concrete form the general rule takes here: no filter may operate on the review's dependent variable.

\paragraph{Row 1: field-restricted LLM block.}

\begin{quote}
\begin{minipage}{0.96\textwidth}
\scriptsize\ttfamily\raggedright
Title:(\\
\quad ``large language model'' OR ``large language models'' OR LLM OR LLMs OR\\
\quad ``generative AI'' OR ``generative artificial intelligence'' OR ChatGPT OR GPT OR\\
\quad Claude OR Gemini OR Llama OR Qwen OR DeepSeek\\
)\\
OR Abstract:(\\
\quad ``large language model'' OR ``large language models'' OR LLM OR LLMs OR\\
\quad ``generative AI'' OR ``generative artificial intelligence'' OR ChatGPT OR GPT OR\\
\quad Claude OR Gemini OR Llama OR Qwen OR DeepSeek\\
)\\
OR Keyword:(\\
\quad ``large language model'' OR ``large language models'' OR LLM OR LLMs OR\\
\quad ``generative AI'' OR ``generative artificial intelligence'' OR ChatGPT OR GPT OR\\
\quad Claude OR Gemini OR Llama OR Qwen OR DeepSeek\\
)
\end{minipage}
\end{quote}

\paragraph{Row 2: broad UML block.}

\begin{quote}
\begin{minipage}{0.96\textwidth}
\scriptsize\ttfamily\raggedright
UML OR ``Unified Modeling Language'' OR PlantUML OR\\
``class diagram'' OR ``class diagrams'' OR ``use case diagram'' OR ``use case diagrams'' OR\\
``sequence diagram'' OR ``sequence diagrams'' OR ``activity diagram'' OR ``activity diagrams'' OR\\
``state machine diagram'' OR ``state machine diagrams'' OR ``state diagram'' OR ``state diagrams'' OR\\
``component diagram'' OR ``component diagrams'' OR ``deployment diagram'' OR ``deployment diagrams'' OR\\
``object diagram'' OR ``object diagrams'' OR ``communication diagram'' OR ``communication diagrams'' OR\\
``package diagram'' OR ``package diagrams'' OR\\
``composite structure diagram'' OR ``composite structure diagrams'' OR\\
``timing diagram'' OR ``timing diagrams'' OR\\
``interaction overview diagram'' OR ``interaction overview diagrams''
\end{minipage}
\end{quote}

\paragraph{Row 3: broad generation/input block.}

\begin{quote}
\begin{minipage}{0.96\textwidth}
\scriptsize\ttfamily\raggedright
generate OR generated OR generating OR generation OR\\
create OR created OR creating OR creation OR\\
construct OR constructed OR construction OR\\
synthesize OR synthesized OR synthesis OR produce OR produced OR\\
transform OR transformed OR transformation OR complete OR completion OR repair OR\\
refine OR refinement OR ``natural language'' OR requirement OR requirements OR\\
specification OR specifications OR ``user story'' OR ``user stories'' OR\\
scenario OR scenarios OR ``problem statement'' OR ``problem statements'' OR\\
``domain description'' OR ``domain descriptions''
\end{minipage}
\end{quote}

\subsection{Manual search}

Manual searching should be employed as a strategy to broaden the scope of the review and minimize the risk of overlooking relevant studies that may not yet be indexed by search engines. The selected data sources are shown in Table~\ref{tab:protocol_manual_venues} and reflects the intersection of Software Modeling and Model Driven Engineering, Requirements Engineering, LLM-based Software Engineering, and empirical evaluation. The selection was pre-defined based on recent secondary studies and venues that contain known seed studies on UML diagram generation by LLMs and validate through consultation with experts in the field\cite{Zhang2027LLMMDE,Rahmanian2026Diagrams,DeBari2024UML,Ferrari2024Sequence,Garaccione2025Class,Garaccione2025UseCase,Xiao2025Sequence}.

\begin{table}[htbp]
\centering
\scriptsize
\caption{Manual venue inspection set}
\label{tab:protocol_manual_venues}
\begin{tabularx}{\textwidth}{p{0.18\textwidth} p{0.37\textwidth} X}
\toprule
\textbf{Status} & \textbf{Venues} & \textbf{Rationale and treatment} \\
\midrule
Core & MODELS main conference; MODELS Companion; MODELSWARD; \emph{Software and Systems Modeling} & Software modeling, UML, and Model Driven Engineering. Proceedings inspected for every review interval. \\
Core & IEEE International Requirements Engineering Conference and related workshops; REFSQ & Relevant to transformation from textual requirements and domain descriptions into software models. \\
Core & ICSE main conference; ICSE Software Engineering in Practice; directly relevant ICSE tracks or workshops; NLBSE & Relevant to LLM-based Software Engineering, industrial evaluation, natural language, and generation tasks. \\
Core & ESEM & Relevant to empirical evaluation, measurement, comparison, and quality assessment. \\
Conditional & EASE; \emph{Empirical Software Engineering} & Activated when the seed set or the automated results show recurrent eligible studies or an indexing gap. \\
Conditional & ECMFA and associated STAF events; SLE & Activated when model transformation, language engineering, or formal modeling studies are recurrent in the candidate corpus. \\
Conditional & ASE; FSE; \emph{Automated Software Engineering}; \emph{Journal of Systems and Software}; \emph{Requirements Engineering}; IEEE TSE; ACM TOSEM & Activated when at least two eligible studies are identified in the outlet, or when a known eligible seed study is not retrieved by the automated search. \\
\bottomrule
\end{tabularx}
\end{table}

A venue is promoted from the conditional set to manual inspection when at least one of the following holds: the seed set contains a relevant study from the venue that no executed automated source retrieves; at least two eligible candidate studies are identified in the venue during automated searching, pilot screening, or citation searching; or the methodological and domain validators jointly determine that excluding the venue creates a documented coverage risk. The venue decision log records the evidence, the years inspected, the route, and the number of unique eligible records contributed. Manual inspection is not expanded because a venue is prestigious or generally relevant to Software Engineering.

\subsection{Consolidated record list}
\label{sec:protocol_record_list}

The automated and manual searches deliver a single consolidated record list. Every record retrieved by any source enters it, one row per record, before any eligibility judgment is made, so that the corpus submitted to screening is a single object with a single provenance history rather than a set of per-database spreadsheets reconciled later.

Three rules govern the list. First, every record receives a stable identifier at entry; the identifier is assigned once, is never reused, and is never renumbered, so that a decision recorded against it remains addressable for the life of the review. Second, every field below is structured, one datum per field, and is transcribed from the source export without editing; when a source supplies a malformed or abbreviated value, the value is preserved as delivered and any correction is recorded beside it rather than written over it. Third, a field that the source does not supply is left empty. A missing metadata value is never inferred, reconstructed from another record, or completed from memory, because a value that was invented is indistinguishable from a value that was retrieved once it sits in the same column.

\begin{longtable}{p{0.20\textwidth} p{0.07\textwidth} p{0.46\textwidth} p{0.14\textwidth}}
\caption{Fields structured in the consolidated record list}
\label{tab:protocol_record_list}\\
\toprule
\textbf{Field} & \textbf{Tag} & \textbf{Content} & \textbf{Obligation} \\
\midrule
\endfirsthead
\toprule
\textbf{Field} & \textbf{Tag} & \textbf{Content} & \textbf{Obligation} \\
\midrule
\endhead
RecordID & --- & Stable review identifier, unique across the whole corpus. & Always \\
SourceDatabase & DB & Database or discovery source that delivered the record. & Always \\
SourceExport & --- & Export file from which the record was read, so that any row can be traced back to a dated retrieval. & Always \\
SourceRecordID & --- & Identifier assigned by the source, when the source assigns one. & When supplied \\
ReferenceType & TY & Record type as declared by the source, such as journal article or conference paper. & Always \\
Title & TI & Title of the report. & Always \\
Authors & AU & Author list, in the order given by the source. & When supplied \\
Editors & A2 & Editors of the containing volume. & When supplied \\
Venue & T2 & Journal, proceedings, or containing volume. & When supplied \\
VenueAbbreviation & J2 & Abbreviated venue name. & When supplied \\
Year & PY & Publication year as declared by the source. & Always \\
Date & DA & Full publication date. & When supplied \\
Volume & VL & Volume. & When supplied \\
Issue & IS & Issue or number. & When supplied \\
StartPage & SP & First page. & When supplied \\
EndPage & EP & Last page. & When supplied \\
DOI & DO & Digital object identifier. & When supplied \\
ISSNorISBN & SN & Serial or book identifier of the containing volume. & When supplied \\
Publisher & PB & Publisher. & When supplied \\
URL & UR & Address at which the record was found. & When supplied \\
Language & LA & Language of the report. & When supplied \\
WorkType & M3 & Type of work as declared by the source, where it differs from the reference type. & When supplied \\
Abstract & AB & Abstract as delivered by the source. & When supplied \\
Keywords & KW & Author or index keywords. & When supplied \\
Affiliation & AD, C1, C3 & Author affiliation and correspondence address, as delivered. Different sources place this in different tags, and each is preserved in the tag that carried it. & When supplied \\
SourceNote & N1 & Note delivered by the source, including contact addresses and publication status remarks. & When supplied \\
\bottomrule
\end{longtable}

The list is the input to deduplication and to the screening procedure. Decision fields, retrieval status, extraction status, and the observations recorded at each gate are added to the same row and are specified in Section~\ref{sec:protocol_records}; they are not part of what the search itself delivers, and the separation is kept so that a re-execution of the search can be compared against this list without decisions from a previous pass contaminating the comparison.

\section{Eligibility criteria}
\label{sec:protocol_eligibility}

Eligibility is defined by two paired lists. Each inclusion criterion is the positive counterpart of exactly one exclusion criterion, and each pair is decided at exactly one point of the screening procedure, called a filter. The \textbf{Filter} column of the tables below names the filter at which each pair is decided; the order of the filters and what is recorded at each of them are given in Section~\ref{sec:protocol_selection}. The criterion states what makes a report eligible, and the procedure states only when the question is asked.

\subsection{Inclusion criteria}

Inclusion codes are numbered in filter order rather than historical order. Each inclusion criterion is the positive counterpart of exactly one exclusion criterion, with the two exceptions stated after the table.

\begin{longtable}{p{0.08\textwidth} p{0.62\textwidth} p{0.08\textwidth} p{0.10\textwidth}}
\caption{Inclusion criteria}
\label{tab:protocol_inclusion}\\
\toprule
\textbf{Code} & \textbf{Criterion} & \textbf{Filter} & \textbf{Pairs with} \\
\midrule
\endfirsthead
\toprule
\textbf{Code} & \textbf{Criterion} & \textbf{Filter} & \textbf{Pairs with} \\
\midrule
\endhead
I1 & The report is within the temporal and language scope of the review. & A1 & E4 \\
I2 & The report is a complete scientific report. & A2 & E1 \\
I3 & The report is a primary study. & A3 & E2 \\
I4 & The study uses at least one LLM as a substantive component of the diagram generation process. & B4 & E6 \\
I5 & The task produces, transforms, completes, repairs, refines, or revises UML diagram content. & B2 & E8 \\
I6 & The source input includes textual software requirements, user stories, scenarios, specifications, problem statements, or explicit textual domain descriptions. & B3 & E9 \\
I7 & The generated result includes explicit UML diagram content, such as PlantUML code, another textual UML representation, XMI, a rendered UML diagram, or another representation from which UML content can be assessed. & B1 & E7 \\
I8 & The UML diagram content is separable, both from other notations and from other generated artifacts, tasks, or outputs. & B5, C1 & E7b, E12 \\
\bottomrule
\end{longtable}

Two structural asymmetries are deliberate. E3 has no positive counterpart, because this criterion predicates a property of the record set rather than of the report. Filter B0 has no criterion of its own, being the conjunction of I4, I5, and I7. As a conjunctive filter, B0 absolves but does not condemn.

I8 must be evaluated at two moments without receiving a separate code for each. At B5, evaluation draws from title and abstract, against mixture with another notation. At C1, evaluation draws from the full text, against mixture with other artifacts, tasks, or outputs. The moment of evaluation must be recorded in the gate outcome field. Creating one code per reading moment would enlarge the table without adding discriminating power.

\subsection{Exclusion criteria}

Exclusion codes are stable identifiers rather than an ordinal sequence. Only the codes listed below are operative, and a code must not be reused or renumbered once screening decisions have been recorded against it, so the numbering is deliberately not contiguous and the gaps carry no meaning. E7b is a sibling of E7 rather than an independent code, because both predicate the same property of the generated result: E7 denotes that the result is not UML, and E7b that it mixes UML with another notation with no separable UML contribution. The suffix marks that relation.

\begin{longtable}{p{0.08\textwidth} p{0.72\textwidth} p{0.08\textwidth}}
\caption{Exclusion criteria}
\label{tab:protocol_exclusion}\\
\toprule
\textbf{Code} & \textbf{Criterion} & \textbf{Filter} \\
\midrule
\endfirsthead
\toprule
\textbf{Code} & \textbf{Criterion} & \textbf{Filter} \\
\midrule
\endhead
E1 & The item is not a complete scientific report, such as an editorial, preface, keynote, tutorial, slide deck, poster-only item, abstract-only record, thesis, book, patent, or non-scientific web content. This list is closed. & A2 \\
E2 & The report is a secondary or tertiary study. & A3 \\
E3 & The report is a duplicate or a less complete member of a publication family. & D \\
E4 & The report is outside the temporal or language scope of the review. & A1 \\
E6 & The study does not use an LLM as a substantive component in producing, transforming, completing, repairing, refining, or revising UML diagram content. & B4 \\
E7 & The generated result is not UML. This includes generic diagram generation and the generation of code, metamodels, or other non-UML artifacts from a UML input. & B1 \\
E7b & The generated result contains UML mixed with another notation, such as C4, ER, BPMN, SysML, Mermaid, or an architecture sketch, with no separable UML contribution. & B5 \\
E8 & The task only evaluates, explains, summarizes, classifies, or discusses an existing UML diagram, without producing, transforming, completing, repairing, refining, or revising UML diagram content. & B2 \\
E9 & The source input is code, an image, an existing UML model, logs, or another non-textual artifact without a substantive textual requirements or domain description component. & B3 \\
E12 & There is no identifiable generation instance: the UML result cannot be separated from other artifacts, tasks, or outputs generated in the same execution, so that no reported result is attributable to the generation of UML content. & C1 \\
\bottomrule
\end{longtable}

If title and abstract are insufficient to decide whether the report involves LLM-based UML generation from textual requirements or domain descriptions, the report must be retained for full-text screening, where the filter that could not be answered is resolved and its code is recorded at that filter. The absence of explicit syntactic or semantic terminology is never sufficient for exclusion, at any gate, because the review's search was deliberately designed without quality terms and selecting on evaluation vocabulary would amount to selecting on the review's own dependent variable.

\subsection{Operational boundary for substantive LLM use}
\label{sec:substantive_llm_use_boundary}

A report satisfies the LLM criterion only when the language model is used as a substantive component in producing, transforming, completing, repairing, refining, or revising the generated UML diagram content. Substantive use means that the model contributes directly to domain-bearing UML content, such as classes, attributes, operations, actors, use cases, relations, messages, lifelines, states, transitions, activities, components, dependencies, or other UML elements and relations.

Model names or broad labels are not sufficient for inclusion. Terms such as pretrained language model, foundation model, or language model support inclusion only when the operational role of the model satisfies the substantive-use requirement. Encoder-only or representation models, including BERT-like models, do not satisfy the LLM criterion solely because they are pretrained language models. When such models are used only as classifiers, sequence taggers, embedding generators, named-entity recognizers, relation extractors, or supervised preprocessing components, the report must be excluded under E6.

Hybrid neural-symbolic pipelines are eligible when the LLM determines, proposes, or revises the semantic content of the UML diagram and symbolic components format, serialize, validate, render, or repair the representation. They are not eligible when symbolic rules determine the UML content and the language model only performs peripheral preprocessing, summarization, formatting, evaluation, or explanation.

When title and abstract do not make the role of the model clear, the report must be retained for full-text screening and filter B4 is deferred, not answered. The exclusion decision must be made only when the full text provides sufficient evidence that the language model does not have a substantive role in producing, transforming, completing, repairing, refining, or revising UML diagram content. A deferred filter remains a filter of its own gate: the decision is read from the full text, and the code and the gate outcome are recorded at B4, never at Gate C.

\subsection{Quality evidence classification during extraction}
\label{sec:protocol_quality_axes}

The review does not define a second eligibility layer, and evidence of quality is not an eligibility filter. All included studies satisfy the same inclusion and exclusion criteria. During data extraction, each included study is classified according to the type and strength of quality evidence that it provides for the generated UML diagram content. These classifications support synthesis without redefining the eligible corpus.

The classification is recorded on three independent axes derived from the operational definitions of Section~\ref{sec:protocol_definitions} and from the model quality relations of Lindland et al. and Krogstie \cite{Lindland1994Quality,Krogstie2006Process}.

\begin{longtable}{p{0.06\textwidth} p{0.16\textwidth} p{0.20\textwidth} p{0.44\textwidth}}
\caption{Quality evidence axes recorded during extraction}
\label{tab:protocol_quality_axes}\\
\toprule
\textbf{Axis} & \textbf{Name} & \textbf{Relation} & \textbf{Values} \\
\midrule
\endfirsthead
\toprule
\textbf{Axis} & \textbf{Name} & \textbf{Relation} & \textbf{Values} \\
\midrule
\endhead
L & Language & Model to language & Absent; textual validity; UML conformity. \\
D & Domain & Model to domain & Absent; claimed only; source requirements; reference model; expert judgment; rubric. \\
U & Use & Model to interpreter and use & Absent; claimed only; comprehension; engineering activity; rework. \\
\bottomrule
\end{longtable}

The three axes are never collapsed into a single scale. Folding pragmatic adequacy into the domain axis would repeat the failure that produced an overloaded exclusion code during the first screening pass, because a fused question yields a fused code, and the distinction becomes unrecoverable.

A fourth field records \textit{result attribution} with values \textit{attributable to UML}, \textit{aggregated with other artifacts}, and \textit{not reported}. It qualifies the strength with which a reported result can be traced to the generation of UML content, in cases where filter C1 confirmed that a generation instance exists.

The extraction additionally classifies which quality constructs the study provides evidence about. That vocabulary is defined once, in the normalized construct field of the quality evidence extraction table, whose syntactic and semantic projections are the two construct fields recorded beside it; it is deliberately not repeated here, because repeating it is what allowed the two enumerations to drift apart. The extraction also records whether syntactic and semantic evidence are reported separately and whether the report provides direct, indirect, absent, unclear, or non-assessable evidence of syntax-semantic dissonance. Axis L is the study-level projection of the syntactic construct field onto the two levels the glossary separates. \textit{Textual validity}, \textit{rendering validity}, \textit{PlantUML parseability}, and \textit{notation-level validity} project onto the textual level, because textual validity is defined as conformance of the concrete output string to the grammar accepted by the selected representation or rendering tool. \textit{UML syntactic conformity}, \textit{metamodel conformance}, and \textit{well-formedness} project onto the UML level, because UML syntactic conformity is defined over the abstract syntax and its well-formedness constraints. No construct label is left without a level, so a report whose only syntactic evidence is a renderer run or a parser run has a value on axis L and is not recorded as absent. Because the axis carries one value per study while the construct field is repeatable, the axis records the most demanding level for which the study reports evidence; this asserts which level was assessed and not that one level implies the other, and the enumeration of what was assessed remains in the construct field. Procedure, reliability, and validity information is recorded in the SQ5 extraction fields and never in the axes.

The analytical subsets are computed from these axes rather than presumed by an eligibility filter.

\begin{longtable}{p{0.24\textwidth} p{0.64\textwidth}}
\caption{Analytical subsets derived from the quality evidence axes}
\label{tab:protocol_analytical_subsets}\\
\toprule
\textbf{Questions} & \textbf{Subset} \\
\midrule
\endfirsthead
\toprule
\textbf{Questions} & \textbf{Subset} \\
\midrule
\endhead
MQ1 to MQ5, and SQ6 & The whole eligible corpus. \\
SQ1 to SQ3 & The whole eligible corpus. \\
SQ4 & Studies for which axis L is not absent and axis D is not absent. \\
SQ7 & Studies for which axis U is not absent. \\
\bottomrule
\end{longtable}

SQ4 asks to what extent the reported evidence distinguishes syntactic from semantic quality. Under an eligibility filter on evidence reporting, that subset would be assumed; under the axes it is computed, which is the difference between describing the literature and describing the review's own filter. The same argument governs MQ5, which asks which combinations remain under-investigated: a gap analysis run over a corpus filtered by evidence reporting would report the filter as a gap in the literature. SQ6 covers the whole corpus for a related reason. What knowledge is supplied as generation context is a property of every generation setting, whereas the pragmatic effects asked by SQ7 exist only where axis U is not absent; reporting the former over the latter subset would describe the supply of domain knowledge and of the UML technical contract over the minority of studies that also measured use. SQ1 to SQ3 cover the whole corpus for the strongest form of the same reason. A quality construct, an inadequacy, and a metric may be purely syntactic: a study whose only evidence is a parser run reports a construct, a definition, and an oracle, and would be silently removed from the construct matrix by a subset conditioned on the domain axis. Restricting the question that carries the constitutive function of this review to the studies that already evaluate against the domain would select on the dependent variable and would report the review's own filter as the vocabulary of the literature.

\subsection{Full-text unavailability as attrition}
\label{sec:protocol_attrition}

A report whose full text is not obtained after the documented access procedure is not excluded. It is placed in a stratum named \textit{identified, not retrieved}. It receives neither an exclusion flag nor an exclusion code, and its gate C outcome is recorded as not retrieved.

Every eligibility criterion predicates something of the report: its year, its type, its object, the operational role of the model, the separability of its result. Full-text unavailability predicates something of the reviewer's access within a time window. Reporting it as an exclusion would read as a statement about the ineligibility of those reports, when the correct statement is that they are reports that could not be read. PRISMA 2020 already separates the two into disjoint boxes, in which \textit{reports not retrieved} precedes and does not overlap \textit{reports excluded, with reasons} \cite{Page2021PRISMA}.

The decision is also empirical. Retrieval loss in this review is systematic rather than random along three dimensions, namely access model, publication type and source, and recency. An exclusion code would hide all three asymmetries, and the recency bias is the most damaging, because the object of the review is a fast-moving literature whose most recent slice is the most informative about model families and prompting strategies.

Reports in the stratum are characterized rather than merely counted, using the metadata and abstracts available for them. The cut-off date for access attempts is 1 September 2026. Its effect is not to exclude reports but to close the stratum: on that date the stratum stops receiving new texts and is counted and described. A text that arrives after the cut-off is recorded with its arrival date and is reported in the sensitivity analysis rather than silently folded into the corpus. Question-level bases are stated explicitly, distinguishing the descriptive questions answerable over all retained reports from the questions that require full text, with a sensitivity check of the former against the abstracts of the non-retrieved stratum.

\subsection{Ambiguity rules}

At title and abstract screening, uncertainty favors retention. A record is excluded only when an exclusion criterion is clearly satisfied. At full-text screening, each decision must cite the relevant criterion and a brief justification. Multi artifact and multi task studies are included when UML specific data can be extracted. Educational studies are eligible because they provide empirical evidence about generation and assessment, but educational context is coded explicitly and considered during interpretation.

\subsection{Publication families}

Reports are grouped into publication families when they share a dataset, artifact, experimental design, tool, or substantial body of results. The most complete report serves as the primary extraction source. Complementary reports are used to fill missing details but are not counted as independent primary studies. Differences between family members are recorded, including changes in models, datasets, methods, or results.

\section{Record management and deduplication}
\label{sec:protocol_records}

All database exports must be preserved in their original formats before transformation. Each record must receive a stable review identifier. Automated deduplication uses DOI, title, authors, year, and source identifiers, followed by manual inspection of probable duplicates. Preprint and published versions are linked rather than deleted before publication family assessment. The master record set must preserves source provenance so that retrieval coverage by database can be analyzed.

The record management log must contain the following fields.

\begin{itemize}
    \item the identification and bibliographic fields of the consolidated record list of Section~\ref{sec:protocol_record_list}, carried unchanged;
    \item duplicate group, publication family identifier, and pre-pass D outcome;
    \item screening decisions by reviewer and stage, recorded per gate as an outcome that names the filter that decided the case;
    \item screening attributes annotated at filters B1 to B4, stored separately from extraction fields and excluded from the agreement calculation;
    \item uncertainty flags raised at title and abstract, which retain the record and never exclude it;
    \item exclusion criterion and justification, with one primary code per excluded record;
    \item full-text retrieval status, including membership of the \textit{identified, not retrieved} stratum;
    \item eligibility decision and quality-evidence classification on axes L, D, and U;
    \item extraction and appraisal status;
    \item recoding flags.
\end{itemize}

\section{Study selection process}
\label{sec:protocol_selection}

Selection comprises calibration, deduplication as a corpus-level pre-pass, title and abstract screening, full-text screening, quality-evidence classification, and publication family reconciliation.

\subsection{Screening procedure}
\label{sec:protocol_gates}

Screening applies the criteria of Section~\ref{sec:protocol_eligibility} as an ordered sequence of single-question decision points, called filters, grouped into gates. This differs from the more common arrangement, in which a flat list of inclusion and exclusion criteria is applied to each report as a whole, and it is adopted for one reason. A flat list records which criterion excluded a report. An ordered sequence of single-question filters also records where the decision was taken and which questions were never reached, so that a second reviewer, or a reviewer working from this protocol in the future, reproduces the route and not only the verdict. Screening is structured so that every decision point carries exactly one question, at most one inclusion criterion, and at most one exclusion criterion. The code records the identity of the criterion; the gate outcome field records where it was applied.

\begin{longtable}{p{0.05\textwidth} p{0.05\textwidth} p{0.46\textwidth} p{0.13\textwidth} p{0.13\textwidth}}
\caption{Screening gate structure}
\label{tab:protocol_gates}\\
\toprule
\textbf{Gate} & \textbf{Filter} & \textbf{Question} & \textbf{Inclusion} & \textbf{Exclusion} \\
\midrule
\endfirsthead
\toprule
\textbf{Gate} & \textbf{Filter} & \textbf{Question} & \textbf{Inclusion} & \textbf{Exclusion} \\
\midrule
\endhead
--- & D & Is the record a duplicate or a less complete member of a publication family? & --- & E3 \\
A & A1 & Is the report within the temporal and language scope? & I1 & E4 \\
A & A2 & Is the report a complete scientific report? & I2 & E1 \\
A & A3 & Is the report a primary study? & I3 & E2 \\
B & B0 & Does the LLM generate or alter the UML content? & I4, I5, I7 & --- \\
B & B1 & Does the generated result include UML content? & I7 & E7 \\
B & B2 & Is the UML content produced or altered? & I5 & E8 \\
B & B3 & Is the source input textual? & I6 & E9 \\
B & B4 & What is the operational role of the model? & I4 & E6 \\
B & B5 & Is the UML content separable from other notations? & I8 & E7b \\
C & C1 & Is there an identifiable generation instance? & I8 & E12 \\
\bottomrule
\end{longtable}

Filter D is a corpus-level pre-pass reported as \textit{records removed before screening}, consistent with PRISMA 2020. Filter A3 routes secondary and tertiary reports to a materialized background pile rather than discarding them. Filter B0 fast-paths clear cases, but a negative answer at B0 does not by itself exclude: it forces descent to B1 through B4 so that the recorded code always originates from a named filter. Filter C1 is the only exclusion available at Gate C.

Filters B1 through B4 are additionally recorded as descriptive attributes even when B0 absolves, without power to exclude. This front-loads three of the review's incommensurability dimensions, namely diagram type, source specification, and model family, at no additional reading cost. Two rules accompany the practice: screening attributes are provisional and are stored in fields separate from extraction fields, and attributes do not enter the inter-rater agreement calculation, which covers only B0 and the named criterion.

Four invariants define the procedure.

\begin{enumerate}[label=(\roman*),leftmargin=*]
    \item \textit{One question per filter.} A filter carries at most one inclusion criterion and at most one exclusion criterion, so no decision is a fused judgment of two properties.
    \item \textit{Fixed order.} Filters are evaluated in the order of Table~\ref{tab:protocol_gates}. When more than one exclusion criterion applies, the first criterion that clearly explains the exclusion in that order is recorded as primary, and the remaining ones are recorded in the justification.
    \item \textit{Stop at the first exclusion.} The first negative answer ends the assessment of that report. The filters that follow are not evaluated and are recorded as not reached.
    \item \textit{Uncertainty retains.} At title and abstract, a report that cannot be decided is retained and carries an uncertainty flag to the next stage. Exclusion requires an exclusion criterion to be clearly satisfied.
\end{enumerate}

The procedure must be executed as follows. Criterion codes are those of Tables~\ref{tab:protocol_inclusion} and~\ref{tab:protocol_exclusion}, whose wording governs; this subsection states execution order and records nothing that those tables do not already define.

\begin{enumerate}[label=\textbf{Step \arabic*.},leftmargin=*,itemsep=0.4em]

    \item \textbf{Pre-pass D, over the record set.} Is the record a duplicate, or a less complete member of a publication family? If yes, remove it before screening, record E3, and stop. This filter predicates a property of the record set rather than of the report, so its output is reported as \textit{records removed before screening} and never as an exclusion with a reason.

    \item \textbf{Gate A, on title and abstract: is this a report the review can read?}
    \begin{enumerate}[label=A\arabic*.,leftmargin=2em]
        \item Is the report within the temporal and language scope? If no, exclude, E4.
        \item Is the report a complete scientific report? If no, exclude, E1.
        \item Is the report a primary study? If no, route it to the background pile, E2. The report is retained as material and is not discarded.
    \end{enumerate}

    \item \textbf{Gate B, on title and abstract: is this the phenomenon the review is about?} Answer B0 first.
    \begin{enumerate}[label=B\arabic*.,leftmargin=2em,start=0]
        \item Does the LLM generate or alter the UML content? If yes, I4, I5 and I7 all hold, the report passes Gate B, and B1 to B4 are still answered and stored as descriptive attributes without power to exclude. If no, or if the answer is not clear, do not exclude here: descend to B1 and answer the filters in order, so that the recorded code always originates from a named filter.
        \item Does the generated result include UML content? If no, exclude, E7.
        \item Is the UML content produced or altered, rather than only evaluated, explained, summarized, classified, or discussed? If no, exclude, E8.
        \item Is the source input textual? If no, exclude, E9.
        \item Is the operational role of the model substantive, as defined in Section~\ref{sec:substantive_llm_use_boundary}? If no, exclude, E6.
        \item Is the UML content separable from other notations? If no, exclude, E7b.
    \end{enumerate}

    \item \textbf{Gate C, on the full text.} Is there an identifiable generation instance, that is, can a reported result be attributed to the generation of UML content rather than to other artifacts, tasks, or outputs produced in the same execution? If no, exclude, E12. C1 is the only exclusion available at Gate C, and it is answered with the full text open.

\end{enumerate}

Two reports illustrate the route. A study that prompts an LLM to produce a class diagram from user stories answers yes at B0 and passes Gate B there, with B1 to B4 annotated as attributes. A study that asks an LLM to explain an existing class diagram answers no at B0, descends, and is excluded at B2 under E8, because it is the filter that names what the study failed to do, and not at B0, which absolves but does not condemn.

Two decisions are deliberately outside this procedure and are never taken at a gate. A report whose full text is not obtained is not excluded: it enters the attrition stratum of Section~\ref{sec:protocol_attrition}. Evidence of quality is not an eligibility condition at any gate: it is classified during extraction on the three axes of Section~\ref{sec:protocol_quality_axes}. Both are stated here because the natural reading of a gate diagram is that everything that reduces the corpus is a gate, and in this review it is not.

\subsection{Reviewer calibration}

Calibration is performed twice, once for each reading moment, because the questions differ. Before production title and abstract screening, the reviewers must independently evaluate a purposive sample containing clearly eligible, clearly ineligible, and ambiguous records, representing different UML diagram types and publication formats. Before full-text screening begins, they must independently calibrate on a purposive sample of full texts covering filter C1 and every filter of Gate A or Gate B liable to be deferred, because whether a generation instance is identifiable is a question that title and abstract calibration never asks. Neither calibration substitutes for the other, and each is dated and recorded separately. Disagreements must be discussed and the eligibility manual revised. Calibration continues until the reviewers achieve at least 80 percent agreement and a Cohen kappa of at least 0.70 for the pilot set \cite{Cohen1960Kappa}. The threshold is a procedural target rather than a claim that kappa alone demonstrates validity.

\subsection{Title and abstract screening}

The primary reviewer screens all deduplicated records. A second reviewer independently screens a stratified random sample of at least 20 percent and every record marked uncertain by the primary reviewer. If disagreement exceeds 10 percent or kappa is below 0.70, an additional 20 percent sample is independently screened. This expansion continues until the predefined threshold is met or all records have been dual screened.

\subsection{Full text screening}

Every record retained after title and abstract screening is assessed independently by two reviewers. Full-text screening performs two distinct operations, and they are recorded in different fields. It resolves any filter of Gate A or Gate B that was deferred for want of evidence, recording the primary exclusion criterion and the gate outcome at the filter that owns that criterion; and it answers filter C1, which is reached only by reports that passed Gate B and whose only exclusion is E12. Each reviewer records one of three decisions, include, exclude, or unclear, and every exclusion carries one primary criterion and a textual justification. Disagreements are resolved through discussion, and unresolved cases are decided by an adjudicator with expertise in Software Engineering and UML or LLM-based modeling.

A concordant \textit{unclear} is not a disagreement and therefore has its own rule. The report is retained, never excluded, and is referred to adjudication with the specific question that could not be answered and the evidence sought. This follows the ambiguity rule of Section~\ref{sec:protocol_eligibility}, under which exclusion requires an exclusion criterion to be clearly satisfied, and it applies with particular force at filter C1, whose question is answered from reported results rather than from a claim. When adjudication cannot resolve it either, the report is retained and the residual uncertainty is carried into extraction as a recorded flag, so that it is visible in the synthesis instead of being settled by default.

Agreement is measured at full text as it is at title and abstract, and it is reported for the two moments separately. It is computed over the primary decision and, for excluded reports, over the filter at which the exclusion was recorded, because agreeing that a report is ineligible while disagreeing about which filter excludes it is a disagreement about the criteria and not about the report. Concordant \textit{unclear} decisions are reported as a count alongside the coefficient and are never absorbed into it. The threshold that governs title and abstract screening governs full-text screening as well; failing it triggers discussion and revision of the eligibility manual rather than an expansion of the sample, because every retained report is already screened by both reviewers.

\subsection{Quality-evidence classification}

Extractable quality evidence is not assessed for eligibility at any gate. For included reports, the type and strength of syntactic, semantic, pragmatic, and correction-related evidence are classified during extraction on the three axes of Section~\ref{sec:protocol_quality_axes}, without creating an additional eligibility status and without affecting the composition of the eligible corpus. A report that evaluates the generated UML content without using the review's terminology is classified, not excluded.

\subsection{Selection reporting}

The final report must be presented in a PRISMA flow diagram and a table of full text exclusions. Counts must distinguish database search, manual search, citation searching, duplicates, non research items, included studies, extraction-derived quality-evidence groupings, and unique primary studies after publication family merging \cite{Page2021PRISMA}.

Two reporting boxes are kept disjoint, in accordance with Section~\ref{sec:protocol_attrition}. \textit{Records removed before screening} carries the corpus-level pre-pass D. \textit{Reports not retrieved} carries the attrition stratum and is reported with its size and with the distribution of its access model, publication type, source, and year, so that the systematic character of the loss is visible rather than absorbed into an exclusion count. \textit{Reports excluded, with reasons} carries only the eligibility criteria of Section~\ref{sec:protocol_eligibility}, each with one primary code and the filter at which it was applied.

\section{Data extraction}
\label{sec:protocol_extraction}

Data extraction is guided directly by the research questions. Two linked sections of one structured extraction form are applied to every included primary study: mapping and generation-context fields, and quality-evidence fields. Each field has a definition, allowed values, examples, missing data code, and relationship to one or more research questions. Analytical subsets are derived from these fields and do not alter eligibility.

\subsection{Mapping and generation-context extraction form}

\begin{longtable}{p{0.17\textwidth} p{0.27\textwidth} p{0.48\textwidth}}
\caption{Mapping and generation-context extraction framework}
\label{tab:protocol_mapping_extraction}\\
\toprule
\textbf{Facet} & \textbf{Field} & \textbf{Extraction rule} \\
\midrule
\endfirsthead
\toprule
\textbf{Facet} & \textbf{Field} & \textbf{Extraction rule} \\
\midrule
\endhead
Bibliographic & Year, venue, publication type, status & Preserve published and online first dates; code peer review and preprint status separately \\
Bibliographic & Authors and affiliations & Extract organizations, countries, and academic or industrial participation \\
Study & Goal and research questions & Preserve verbatim formulation and code study purpose \\
Study & Research method & Classify experiment, quasi experiment, benchmark, case study, user study, demonstration, design and evaluation, qualitative study, or other \\
Study & Contribution type & Tool, method, model, framework, dataset, benchmark, empirical findings, taxonomy, or other \\
Artifact & DiagramType & Class diagram; use case diagram; sequence diagram; activity diagram; state machine diagram; state diagram; component diagram; deployment diagram; object diagram; communication diagram; package diagram; composite structure diagram; timing diagram; interaction overview diagram; multiple UML types; unclear. \\
Artifact & GenerationTask & Generation; transformation; completion; repair; refinement; iterative revision; mixed task; unclear. \\
Input & InputRequirementForm & Requirements document; user stories; use case narratives; scenarios; problem statements; domain descriptions; structured specifications; mixed textual input; unclear. \\
Input & Input structure and quality & Free text, template, semi structured, structured; reported ambiguity, inconsistency, incompleteness, or preprocessing \\
Domain and data & Application domain and dataset & Domain, public or private status, size, provenance, reference availability \\
LLM & ModelFamily & Exact model family or model name reported by the study, such as GPT, ChatGPT, Claude, Gemini, Llama, Qwen, DeepSeek, T5, BART, BERT, RoBERTa, or other. \\
LLM & ModelArchitectureOrAccessMode & Autoregressive; encoder-decoder; encoder-only; conversational/instructional; API-based; open-weight; fine-tuned; unclear; not reported. \\
LLM & GenerativeOrInstructionalUse & Yes; No; Unclear. \\
LLM & LLMRoleInPipeline & Semantic content generator; semantic content reviser; candidate proposer; surface formatter; syntax repairer; validator; evaluator only; preprocessor only; embedding or classifier only; unclear. \\
LLM & SemanticContentAuthority & LLM-determined; hybrid LLM-symbolic; symbolic/rule-determined; unclear. \\
LLM & SymbolicPipelineRole & None reported; formatting; serialization; validation; rendering; syntax repair; semantic rule generation; content generation; unclear. \\
LLM & SubstantiveLLMUseDecision & Included as substantive LLM use; excluded as non-substantive language-model use; unclear at title and abstract; resolved at full text. \\
LLM & BoundaryDecisionRationale & Short textual explanation of how the reviewers decided whether the model had a substantive role in producing, transforming, completing, repairing, refining, or revising UML diagram content. \\
Context & DomainKnowledgeProvided & None reported; implicit in requirements; explicit domain description; examples; ontology or structured knowledge; retrieval augmented context; expert-provided context; unclear. \\
Context & UMLTechnicalInstructionProvided & None reported; notation instruction; metamodel or well-formedness instruction; PlantUML syntax instruction; examples; tool or parser feedback; unclear. \\
Context & PromptingOrContextStrategy & Zero-shot; few-shot; chain-of-thought or reasoning instruction; retrieval augmented generation; iterative prompting; agentic workflow; fine tuning; tool feedback loop; mixed; unclear. \\
Context & Adaptation and orchestration & Fine tuning, preference optimization, agents, tools, self correction, iterative feedback, or none \\
Inference & Sampling and repetition & Temperature, top p, top k, seed, number of generations, selection policy \\
Output & OutputRepresentation & PlantUML code; rendered UML diagram; XMI; textual UML description; image; mixed representation; other explicit UML representation; unclear. \\
Output & PlantUMLGenerated & Yes; No; Unclear. \\
Output & NotationFamily & UML only; UML with a standard profile; a UML-derived language that is not a UML profile; UML alongside a separable non-UML notation; unclear. Records against which definition of the language syntactic conformity was assessed. Filter B5 establishes that a separable UML contribution exists and then discards which other notation it was; this field keeps what the filter discards. \\
Evaluation & BaselineCondition & No baseline; human-authored model; rule-based or deterministic tool; another LLM; another version or configuration of the same LLM; ablation of the generation context; published benchmark or gold standard; other. Repeatable. Reported quality dimensions are not recorded here: they are given by the three axes and by the dimension partition of the normalized construct field. \\
Open science & Data, prompts, code, models, and replication package & Record availability, persistent identifier, access restrictions, and completeness \\
\bottomrule
\end{longtable}

\subsection{Quality-evidence extraction form}

\begin{longtable}{p{0.17\textwidth} p{0.27\textwidth} p{0.48\textwidth}}
\caption{Quality-evidence extraction framework}
\label{tab:protocol_quality_evidence_extraction}\\
\toprule
\textbf{Facet} & \textbf{Field} & \textbf{Extraction rule} \\
\midrule
\endfirsthead
\toprule
\textbf{Facet} & \textbf{Field} & \textbf{Extraction rule} \\
\midrule
\endhead
Axis & QualityAxisL, QualityAxisD, QualityAxisU & One field per axis, whose values are those of Table~\ref{tab:protocol_quality_axes} and are not restated here. The same rule that keeps the construct vocabulary in a single place applies to the axes: an enumeration written twice is an enumeration that drifts. \\
Axis & ResultAttribution & Attributable to UML; aggregated with other artifacts; not reported. Qualifies how far a reported result can be traced to the generation of UML content. \\
Construct & Original quality term & Preserve the term exactly as reported, before any mapping to the review's scheme \\
Construct & Definition and theoretical source & Extract verbatim or near verbatim definition within copyright limits; record cited framework \\
Construct & Normalized construct & Single normalization vocabulary of quality \textit{attributes}, partitioned by dimension. Syntactic: textual validity; rendering validity; PlantUML parseability; UML syntactic conformity; metamodel conformance; well-formedness; notation-level validity. Semantic: semantic correctness; semantic completeness; requirements coverage; domain fidelity; consistency; consistency with source requirements. Pragmatic: pragmatic adequacy; understandability; readability; clarity. Or other. Deviation labels are not admitted here and are recorded as a discrepancy operation and a carrier in the inadequacy fields below; effort quantities are not admitted here and are recorded in the rework fields. Section~\ref{sec:protocol_attribute_deviation_measure} states the duality. The two construct fields below are the syntactic and semantic projections of this vocabulary and admit no value absent from it. \\
Construct & SyntacticEvidenceAvailable & Yes; No; Unclear. \\
Construct & SemanticEvidenceAvailable & Yes; No; Unclear. \\
Construct & SyntacticQualityConstruct & Textual validity; rendering validity; PlantUML parseability; UML syntactic conformity; metamodel conformance; well-formedness; notation-level validity; other; not reported. \\
Construct & SemanticQualityConstruct & Semantic correctness; semantic completeness; requirements coverage; domain fidelity; consistency; consistency with source requirements; other; not reported. \\
Metric & SyntacticMetricOrProcedure & Exact metric, rubric, parser, compiler, rendering check, metamodel check, checklist, human judgment, automated tool, or narrative procedure reported by the study. \\
Metric & SemanticMetricOrProcedure & Exact metric, rubric, oracle, reference model, gold standard, requirements traceability procedure, expert judgment, checklist, human coding, automated comparison, or narrative procedure reported by the study. \\
Reference & EvaluationReference & Requirements text; domain description; reference UML model; gold standard; rubric; expert judgment; benchmark; parser or renderer; UML specification; metamodel; other; unclear. \\
Reference & OracleOrGoldStandard & Yes; No; Partial; Unclear. \\
Reference & Reference construction and validation & Who created it, from which sources, through which process, and whether disagreement or alternative valid models were handled \\
Unit & Unit of assessment & Diagram, requirement, proposition, actor, use case, class, attribute, relation, message, state, transition, or other \\
Inadequacy & Original label and definition & Preserve the source label and its operational meaning \\
Inadequacy & Normalized discrepancy operation & Omission, unsupported addition, incorrect representation, incorrect classification, inconsistency, duplication, redundancy, or emergent category. These are \textit{deviations}, the dual of the normalized construct field above and never a repetition of it; Section~\ref{sec:protocol_attribute_deviation_measure} states the correspondence. Selecting \textit{emergent category} requires recording the emergent label in free text \\
Inadequacy & Violated reference & Concrete representation, UML abstract syntax, UML construct semantics, domain or requirements semantics, cross diagram consistency, or intended engineering task \\
Inadequacy & UML carrier & Affected diagram, element, relation, property, constraint, sequence, state, transition, message, behavior, or other type specific unit \\
Inadequacy & Reported severity or task effect & Severity scale or task effect as the study reports it, in the study's own terms. This field exists because Equation~\ref{eq:protocol_inadequacy_tuple} makes severity the fourth dimension of the coding tuple; it is left empty when the study reports none, and emptiness is not imputed \\
Metric & Metric name and formula & Extract formula, numerator, denominator, scale, direction, threshold, and aggregation rule \\
Metric & Automation level & Fully automated, tool assisted, human coded, LLM assessed, or mixed \\
Evaluation & HumanEvaluatorRole & Author; independent evaluator; domain expert; UML expert; student; practitioner; crowd worker; participant; not reported; not applicable. Repeatable, one occurrence per role. \\
Evaluation & AutomatedEvaluatorTool & PlantUML; parser; renderer; model checker; custom script; LLM-as-judge; other; none reported. \\
Evaluation & QualityAggregationRule & Per diagram; per element; per relation; per requirement; per prompt; per model; percentage; score; ordinal scale; pass/fail; narrative only; unclear. \\
Reliability & Agreement and uncertainty & Agreement measure, value, confidence interval, repeated generation, variance, effect size, significance, or qualitative confidence \\
Context relation & Context factor & Factor manipulated, controlled, observed, or merely reported \\
Context relation & Inferential status & Descriptive, comparative, relational, causal, mechanistic, or unsupported claim \\
Pragmatic use & Task, actor, and outcome & Engineering activity, participant profile, success criterion, time, accuracy, decision quality, or comprehension \\
Rework & CorrectionOrReworkEvidence & Yes; No; Unclear. \\
Rework & Correction measure & Time, edit count, changed elements, distance, severity, number of iterations, subjective effort, or other \\
Syntax-semantic dissonance & SyntaxSemanticDissonanceEvidence & Direct; Indirect; Absent; Unclear; Not assessable. \\
Syntax-semantic dissonance & SyntaxSemanticDissonanceType & Syntactically valid but semantically incorrect; syntactically valid but semantically incomplete; renderable but domain-inconsistent; textually valid PlantUML with wrong UML relation; syntactically invalid but semantically partially recoverable; other; not applicable. \\
Syntax-semantic dissonance & DissonanceEvidenceBasis & Short textual explanation of the evidence supporting the classification. \\
Limitations & Construct and evaluation threats & Author reported and reviewer inferred limitations, kept in separate fields \\
\bottomrule
\end{longtable}

\subsection{Missing and ambiguous data}

The extraction database distinguishes \textit{not reported}, \textit{not applicable}, \textit{unclear}, and \textit{not accessible}. Reviewers must not infer missing model versions, prompts, parameters, metrics, or evaluator qualifications from external assumptions. Supplementary material and replication packages may be used when explicitly linked to the report. Authors may be contacted once for clarification when the missing information is essential to quality-evidence interpretation, and every contact attempt is logged.

\subsection{Extraction pilot and verification}

The forms must be piloted on a temporally and methodologically diverse sample of at least ten studies. The sample must include different UML diagram types and evaluation references. The pilot assesses field clarity, extraction time, category overlap, missing data frequency, and the feasibility of unit level extraction. All interpretive quality-evidence fields are independently extracted by two reviewers. Objective mapping metadata may be extracted by one reviewer and verified through an independent audit of at least 20 percent. Any material revision to the form or codebook triggers retrospective review of studies already extracted.

\section{Methodological and reporting appraisal}
\label{sec:protocol_appraisal}

Included studies are not excluded solely because of methodological weakness. Exclusion based on a methodological quality score could distort the map of an emerging area. All included studies are appraised to qualify the credibility and comparability of their quality evidence. Appraisal distinguishes methodological adequacy, reporting completeness, and availability of replication material. No single total score will be used as a substitute for domain specific judgment.

\begin{longtable}{p{0.10\textwidth} p{0.58\textwidth} p{0.24\textwidth}}
\caption{Methodological and reporting appraisal domains}
\label{tab:protocol_appraisal}\\
\toprule
\textbf{Code} & \textbf{Appraisal question} & \textbf{Judgment} \\
\midrule
\endfirsthead
\toprule
\textbf{Code} & \textbf{Appraisal question} & \textbf{Judgment} \\
\midrule
\endhead
A1 & Are the study objective and research questions explicit and consistent with the evaluation? & Yes, Partially, No, Not applicable \\
A2 & Are the requirements, domain material, dataset provenance, sample size, and selection procedure described? & Yes, Partially, No, Not applicable \\
A3 & Are the LLM identity, version, access mode, instructions, context, tools, parameters, and execution date sufficiently reported? & Yes, Partially, No, Not applicable \\
A4 & Is the UML diagram type and output representation clearly specified? & Yes, Partially, No, Not applicable \\
A5 & Are the quality constructs defined and operationalized? & Yes, Partially, No, Not applicable \\
A6 & Is the evaluation reference explicit, inspectable, and justified? & Yes, Partially, No, Not applicable \\
A7 & Are metrics, rubrics, units, thresholds, and aggregation rules sufficiently defined? & Yes, Partially, No, Not applicable \\
A8 & Are evaluator qualifications, training, independence, and disagreement procedures reported when human judgment is used? & Yes, Partially, No, Not applicable \\
A9 & Is stochastic variation addressed through repeated generation, sampling control, or uncertainty reporting? & Yes, Partially, No, Not applicable \\
A10 & Are baselines and comparison conditions appropriate to the claim? & Yes, Partially, No, Not applicable \\
A11 & Are quantitative or qualitative analyses appropriate and traceable to the data? & Yes, Partially, No, Not applicable \\
A12 & Are threats, limitations, missing data, and failed cases reported? & Yes, Partially, No, Not applicable \\
A13 & Are prompts, outputs, datasets, code, rubrics, or replication materials available? & Yes, Partially, No, Not applicable \\
\bottomrule
\end{longtable}

Four domains are treated as critical for interpreting quality-evidence findings: A3, A5, A6, and A7. Sensitivity analyses compare conclusions derived from all included studies with conclusions derived from studies that have no \textit{No} judgment in the critical domains. The appraisal is not used to manufacture a universal quality rank among heterogeneous study designs.

\section{Coding and taxonomy construction}
\label{sec:protocol_coding}

The taxonomy must be constructed through a hybrid deductive and inductive process. The upper level structure is deductively informed by conceptual modeling quality and the UML specification \cite{Lindland1994Quality,Krogstie2006Process,OMG2017UML}. Diagram specific categories and terminology are developed inductively from primary studies through thematic synthesis and constant comparison \cite{Cruzes2011Thematic,Krippendorff2018Content}.

\subsection{Coding dimensions}

Every extractable inadequacy is represented, where evidence permits, as a tuple

\begin{equation}
I = \langle R_v, O_d, C_u, S_e, E_x \rangle,
\label{eq:protocol_inadequacy_tuple}
\end{equation}

where $R_v$ is the violated reference, $O_d$ is the discrepancy operation, $C_u$ is the UML carrier, $S_e$ is the reported or inferred severity, and $E_x$ is the example or evidence excerpt.

The initial dimensions are the following.

\begin{enumerate}[label=\roman*)]
    \item \textbf{violated reference}, comprising concrete representation, UML abstract syntax, UML construct semantics, domain or requirements semantics, cross diagram consistency, or intended engineering task;
    \item \textbf{discrepancy operation}, comprising omission, unsupported addition, incorrect representation, incorrect classification, inconsistency, duplication, redundancy, or emergent category;
    \item \textbf{UML carrier}, comprising the affected diagram, element, relation, property, constraint, sequence, state, transition, message, behavior, or other type specific unit;
    \item \textbf{severity or consequence}, preserved when the study reports a severity scale or task effect;
    \item \textbf{provenance}, comprising the original term, definition, page, table, figure, metric, or example supporting the code.
\end{enumerate}

\subsection{Coding procedure}

\begin{enumerate}[label=\arabic*)]
    \item Extract original terminology and definitions without immediately replacing them with normalized labels.
    \item Apply the initial deductive dimensions to a pilot subset.
    \item Create open codes for categories that cannot be represented without loss.
    \item Compare codes across studies and diagram types, merge synonymous codes, separate overloaded codes, and document boundaries with positive and negative examples.
    \item Construct diagram specific subcategories under shared transversal dimensions.
    \item Search for negative cases that challenge category definitions.
    \item Independently recode the pilot subset with the revised codebook.
    \item Apply the stable codebook to the remaining included corpus.
    \item Retrospectively recode earlier studies whenever a material category or boundary changes.
\end{enumerate}

The codebook must be versioned. Each code contains a definition, inclusion rule, exclusion rule, positive example, negative example, parent category, allowed UML carriers, and revision history. Original study labels remain available in a separate field so that normalization is fully auditable.

\subsection{Coding reliability}

Coding an inadequacy is two operations, and only the second carries an agreement coefficient. The first is unitization. The inadequacies reported by a study are enumerated in the study's own vocabulary before any normalized label is applied, and two reviewers may enumerate different numbers of them from the same study. The enumerated lists are therefore reconciled by discussion and adjudication before classification begins, and the reconciliation is recorded as the count of items added, removed, and merged on each side. Enumeration disagreement is reported as that count and is never absorbed into a coefficient computed over classification.

The second operation is classification. Two reviewers independently classify every item of the reconciled list, for every included study. Agreement is assessed at the normalized category level and at the principal dimensions of the tuple in Equation~\ref{eq:protocol_inadequacy_tuple}. Cohen kappa is used for binary and mutually exclusive nominal decisions \cite{Cohen1960Kappa}. Krippendorff alpha is preferred when categories are multiple, data are missing, or more than two coders are involved \cite{Krippendorff2018Content}. Reliability values are interpreted together with confusion matrices, category prevalence, and documented disagreement patterns. Disagreements are resolved by discussion and adjudication, and the consensus code is stored separately from the original reviewer codes.

Reconciling the units before classifying them is what makes a contingency table exist. An inadequacy is a repeatable field of free enumeration, so without a shared list of items there is no pairing between the two reviewers' codes, and a coefficient computed over unpaired enumerations would report how much each reviewer found rather than how far they agree on category boundaries. The rule is fixed here, before the two independent passes, because aligning the units afterwards would let each reviewer's classification decide which items are held to have been reported at all.

\section{Synthesis plan}
\label{sec:protocol_synthesis}

\subsection{Descriptive mapping of included studies}

The included studies are mapped through frequencies, proportions, temporal distributions, cross tabulations, and co-occurrence matrices. Planned analyses include the following.

\begin{itemize}
    \item publication year by venue and publication type;
    \item diagram type by generation task;
    \item diagram type by input source and output representation;
    \item model family and version by year;
    \item context strategy by diagram type and task;
    \item study design by contribution type;
    \item quality dimension assessed, as given by the three evidence axes and by the dimension partition of the normalized construct, by diagram type;
    \item evaluation measure and procedure by generation task;
    \item baseline and human involvement by study design;
    \item baseline condition by generation task and diagram type;
    \item replication material availability by year and venue;
    \item publication status and methodological characteristics.
\end{itemize}

The co-occurrence matrices are computed from the cross tabulations listed above, and three of them carry the concentration analysis: diagram type by generation task, diagram type by evaluation measure, and generation configuration by evaluation measure, where the generation configuration is the pair of model family and context strategy. A cell is reported as under investigated when fewer than two studies populate it. The threshold is fixed before the matrices are computed and is reported with them, because a threshold chosen after the distribution is known would select the gaps that it then reports. A cell that no study could populate is distinguished from a cell that no study did populate, because a combination that the corpus makes impossible is not a gap in the literature.

Counts are reported using both reports and unique primary studies where relevant. Categories with multiple values are explicitly identified as non exclusive. No meta analysis is planned because heterogeneous tasks, units, references, and metrics are expected. Any later quantitative aggregation requires a documented decision and a demonstrated basis for commensurability.

\subsection{Extraction-based synthesis of quality constructs}

Definitions are synthesized through a construct matrix containing the original term, reported definition, theoretical source, normalized construct, evaluation reference, unit, indicator, metric, and study context. The synthesis will identify convergent definitions, overloaded terms, unnamed operationalizations, and cases in which identical labels denote different procedures. The synthesis reports the number of distinct original terms and the number of normalized constructs into which they collapse. A ratio close to one to one would show that the normalization did no work, and is reported as such rather than presented as a vocabulary of the literature.

\subsection{Taxonomy of reported UML inadequacies}

The taxonomy must be presented at two levels. The transversal level organizes discrepancy operations and violated references that recur across diagram types. The specific level records how these operations manifest in actors, use cases, classes, relations, attributes, messages, lifelines, states, transitions, activities, and other UML elements. Frequencies are presented as frequencies of reporting and coding within the reviewed corpus, not as estimates of real world defect prevalence unless the primary data support that inference. The taxonomy reports the number of distinct original labels and the number of normalized discrepancy operations they collapse into, on the same terms.

\subsection{Metric and evaluation-reference catalogue}

Metrics and procedures are organized by the measurement tail of the spine of Equation~\ref{eq:protocol_quality_spine}, defined in Section~\ref{sec:protocol_attribute_deviation_measure}, namely the chain

\begin{equation}
\text{construct} \rightarrow \text{reference} \rightarrow \text{unit} \rightarrow \text{indicator} \rightarrow \text{procedure} \rightarrow \text{aggregation}.
\label{eq:protocol_metrology_chain}
\end{equation}

The synthesis compares parsability, compilation, UML rule checks, precision, recall, F scores, coverage, edit distance, graph similarity, expert ratings, LLM judge scores, task performance, time, and other measures only after reconstructing their definitions and denominators. Measures with the same name but different operationalizations are not pooled.

\subsection{Syntax-semantic dissonance synthesis}

Syntax-semantic dissonance is analyzed from the explicit extraction fields for syntactic evidence, semantic evidence, dissonance classification, dissonance type, and evidence basis. The synthesis distinguishes direct evidence from indirect evidence and does not infer semantic adequacy from parseability, renderability, textual validity, or UML syntactic conformity alone.

Analytical subsets include studies with syntactic evidence only; studies with semantic evidence only; studies with both syntactic and semantic evidence; studies with direct syntax-semantic dissonance evidence; studies with indirect or non-assessable dissonance evidence; and studies generating PlantUML code. These subsets are derived during extraction and are not eligibility layers.

\subsection{Generation context knowledge synthesis}

This synthesis covers the whole eligible corpus, because what knowledge is supplied as generation context is a property of every generation setting and not only of the studies that also measured use. The synthesis examines how requirements, technical specifications, domain knowledge, UML technical instructions, prompts, examples, retrieval, tools, parser or renderer feedback, and iterative workflows are provided as generation context, and reports the supply of domain knowledge and the supply of the UML technical contract as the two axes of the resulting map.

\subsection{Pragmatic adequacy and rework synthesis}

This synthesis covers only the studies for which axis U is not absent, because a pragmatic or rework effect exists only where use was assessed, and running it over the whole corpus would read the silence of the studies that never measured use as an absence of effect. Reported context factors are classified as merely described, experimentally manipulated, statistically associated, or proposed without empirical evaluation, and causal language is used only when the original design supports it. Pragmatic and rework evidence is synthesized by task, actor, diagram type, outcome, and measure. A study that reports comprehensibility without an engineering task is not automatically coded as evidence of pragmatic adequacy.

\subsection{Heterogeneity and sensitivity analysis}

Heterogeneity is examined across diagram type, educational or professional context, input type, dataset provenance, model access, model family, publication status, evaluation reference, and appraisal domains. Sensitivity analyses will compare, where the corpus permits, peer reviewed versus preprint studies, expert reference versus LLM judge evaluation, repeated versus single generation studies, and stronger versus weaker reporting in critical appraisal domains.

\section{Protocol validation}
\label{sec:protocol_validation}

The protocol is validated before production screening and before extraction through five complementary procedures. Validation is a precondition for execution rather than a report of it.

\begin{enumerate}[label=\arabic*)]
    \item \textbf{Conceptual validation}. The scope, the operational definitions, the questions, and the extraction fields are reviewed by the domain validators of Table~\ref{tab:protocol_reviewers}.
    \item \textbf{Search validation}. The search strategy is tested against the seed set, inspected by a methodological validator, and revised when a failure reveals an avoidable retrieval gap.
    \item \textbf{Selection pilot}. The reviewers apply the filters of Table~\ref{tab:protocol_gates} to a heterogeneous sample, and ambiguous criteria and examples are revised. Calibration must meet the target of Section~\ref{sec:protocol_gates} before production screening.
    \item \textbf{Extraction pilot}. The extraction form is applied to at least ten studies and revised for clarity, feasibility, and coverage.
    \item \textbf{Synthesis pilot}. The initial coding architecture is applied to a subset covering more than one diagram type and more than one evaluation procedure, and the categories are revised before full coding.
\end{enumerate}

Two rules govern the independence of validation, and they exist because the review team is small enough that one researcher can hold more than one role. First, validation is completed and dated before the activity it validates begins, so that no rule is approved after its consequences are visible. Second, when a disagreement referred to adjudication turns on a protocol rule, the adjudication record names the validator who approved that rule, so that the concentration is visible in the audit trail rather than absorbed into a verdict. The UML domain validator holds no adjudication role and is therefore the independent check on the codebook boundaries, the UML carrier categories, and the reference distinctions.

Validation results, comments, decisions, and the resulting protocol changes are stored in the validation record together with the version of the protocol to which they refer. Protocol changes before execution are recorded in the version changelog. Formal amendments in the sense of PRISMA-P are reopened when execution begins, and the decision to defer them is itself recorded \cite{Shamseer2015PRISMAP}.

\section{Threats to validity and mitigation}
\label{sec:protocol_validity}

\begin{longtable}{p{0.18\textwidth} p{0.34\textwidth} p{0.40\textwidth}}
\caption{Principal threats and mitigation strategies}
\label{tab:protocol_threats}\\
\toprule
\textbf{Threat} & \textbf{Potential consequence} & \textbf{Mitigation} \\
\midrule
\endfirsthead
\toprule
\textbf{Threat} & \textbf{Potential consequence} & \textbf{Mitigation} \\
\midrule
\endhead
Search terminology & Relevant studies may omit the LLM, UML, generation, or requirements terminology used in the query & Broad synonym blocks, model names, diagram type names, seed validation, manual venue search, and snowballing \\
Database coverage & Recent workshop or preprint studies may be missing from the major indexes & Three executed indexes, direct digital libraries, citation searching, manual venue search, and a final search update \\
Unactivated complementary sources & Web of Science, ScienceDirect, and arXiv are planned and not executed, so a study indexed only by one of them is absent from the corpus rather than excluded by a criterion & The gap is declared rather than presented as a decision, the sources carry an explicit pending status in Table~\ref{tab:protocol_sources}, and the source contribution analysis reports it. Activation or a justified exclusion is required before the search is declared closed \\
Venue selection bias & A predefined manual list may overrepresent the expected Software Engineering communities & Automated sources define the principal candidate corpus and manual venues complement it \\
Scope ambiguity & Generic diagrams or non-requirements inputs may be inconsistently included & Operational definitions, examples, calibration, dual full-text screening, and adjudication \\
Publication family duplication & Conference, preprint, and journal versions may be counted as independent evidence & Family identification, primary report selection, complementary extraction, and unique study counts \\
Retrieval loss & Full texts not obtained are systematically distributed by access model, publication type, and recency, and the most recent slice is the most informative & Attrition stratum rather than an exclusion code, characterization of the stratum from metadata and abstracts, a stated cut-off date, and question-level bases with a sensitivity check \\
Construct heterogeneity & Identical labels may denote different quality constructs or metrics & Original term preserved before normalization, construct matrix, reference and unit extraction, and no pooling without commensurability \\
Pole confusion & A deviation type may be recorded as a quality attribute, so that the construct matrix reports defect types as dimensions of the literature & The two normalization vocabularies are disjoint duals, stated in Section~\ref{sec:protocol_attribute_deviation_measure} and enforced by a lock over the codebook \\
Stipulated terminology & \textit{Inadequacy} is stipulated rather than borrowed, and a reader may import the pragmatic or the natural language generation sense & The stipulation names both excluded readings, the declined standard alternatives are justified, and each deviation is additionally named by the axis of the reference it violates \\
Taxonomic normalization & Reviewer categories may overwrite source meanings & Separate original and normalized fields, codebook provenance, negative cases, double coding, and retrospective recoding \\
Unit misalignment & Two reviewers may enumerate different numbers of inadequacies from one study, so that an agreement coefficient would measure how much each found rather than how far they agree & Unitization is reconciled before classification, the reconciliation is reported as a count, and only classification carries a coefficient \\
Reviewer role concentration & One researcher may hold adjudication and validation roles, so that a rule is adjudicated by the person who approved it & Validation is dated and completed before the activity it validates; adjudication records which validator approved the rule in dispute; the UML domain validator holds no adjudication role \\
Reviewer subjectivity & Eligibility, extraction, appraisal, or coding may vary by reviewer & Training, calibration, independent decisions, agreement analysis, adjudication, and a retained audit trail \\
Reporting bias & Studies may omit unsuccessful generations, inadequacies, parameters, or evaluator disagreement & Appraisal of reporting, missing data coding, artifact inspection, author contact when essential, and cautious inference \\
Rapid model evolution & Findings may become technologically dated & Exact model and version extraction, execution date coding, temporal analysis, and a search update \\
Preprint inclusion & Non peer reviewed evidence may distort conclusions & Explicit status coding, separate analysis, and sensitivity analysis \\
LLM assisted review work & Automated assistance may introduce untraceable or fabricated decisions & Human responsibility for every inclusion, extraction, and code, and logging of tools, outputs, and verification \\
Selective protocol change & Criteria or outcomes may be changed after results are observed & Every pre-execution change is recorded in the version changelog with its rationale, executable locks assert the invariants that the prose states, and formal amendments are reopened at the start of execution \\
\bottomrule
\end{longtable}

\section{Data management, transparency, and use of artificial intelligence}
\label{sec:protocol_data}

The replication package will contain the registered protocol, search strings, search dates, raw exports when licensing permits, manual-search source manifests, permissible source snapshots or reproducibility hashes, raw documentary inventories, normalized metadata, inventory-conflict records, screening decisions, full text exclusion reasons, publication family links, extraction forms, codebook versions, appraisal judgments, analysis scripts, figures, and tables. Licensed database exports, publisher pages, and copyrighted full texts are redistributed only when their terms permit it; otherwise the package preserves the query, source URL, retrieval date, permitted audit metadata, and checksums or reproduction instructions sufficient to document the retrieval without republishing restricted content \cite{Rethlefsen2021PRISMAS}.

LLM or automation tools may be used for clerical support such as formatting, candidate duplicate detection, translation of search syntax, or consistency checking. They will not make autonomous eligibility, extraction, appraisal, or taxonomy decisions. Every substantive decision remains attributable to a human reviewer. Tool name, version, purpose, prompt or rule, output, and verification procedure are logged. Model generated bibliographic data are verified against authoritative records before use.

The review uses published literature and does not initially involve human participants. If authors are contacted for clarification, correspondence is limited to study information and stored according to institutional data governance requirements. Reviewer conflicts of interest are documented.

\section{Reporting and completion criteria}
\label{sec:protocol_reporting}

The secondary study will be considered methodologically complete when the following conditions are satisfied.

\begin{enumerate}[label=\arabic*)]
    \item The validated search has been executed in all retained automated sources and updated before final reporting.
    \item The complete core manual venue set has been inspected, every conditional venue trigger has been resolved, and the incremental yield of manual searching has been reported.
    \item Every required manual-search unit has a documented operational state; before final synthesis, every available unit has reached \texttt{COMPLETE}, while genuinely unavailable future proceedings remain explicitly \texttt{PENDING}. Source manifests, raw and normalized inventories, discovery logs, candidate sets, and coverage comparisons are preserved in the replication package.
    \item Source contribution, redundancy, seed recovery, and any source retention or removal decision are documented.
    \item Deduplication, selection, and publication family reconciliation are complete and auditable.
    \item Every included study has a verified mapping and generation-context extraction record.
    \item Every included study has independently verified quality-evidence extraction and appraisal.
    \item The codebook has reached a stable version and all included studies have been coded under that version.
    \item Research question specific syntheses have been completed and checked against the extracted data.
    \item Sensitivity analyses and negative case checks have been performed where applicable.
    \item The PRISMA flow, SEGRESS checklist, PRISMA-S search record, exclusion table, and replication package are complete.
    \item Conclusions distinguish the mapped corpus, extraction-derived analytical subsets, publication reports, unique primary studies, and evaluation instances.
\end{enumerate}

The final report will avoid causal or prevalence claims that exceed the designs and units represented in the primary studies. Frequencies in the map describe the literature, not necessarily the real frequency of inadequacies in LLM generated UML diagrams. The taxonomy describes reported and operationalized inadequacies, not a complete ontology of every defect that can occur.

\section{Protocol validation checklist}
\label{sec:protocol_checklist}

The checklist is a precondition for execution. Items V01 to V09 must be complete before production screening begins, and items V10 to V14 before extraction begins.

\begin{longtable}{p{0.06\textwidth} p{0.70\textwidth} p{0.16\textwidth}}
\caption{Checklist for protocol validation before execution}
\label{tab:protocol_checklist}\\
\toprule
\textbf{ID} & \textbf{Validation item} & \textbf{Status} \\
\midrule
\endfirsthead
\toprule
\textbf{ID} & \textbf{Validation item} & \textbf{Status} \\
\midrule
\endhead
V01 & The mapping design and the quality-evidence emphasis are explicitly justified. & \protocolfield{Pending} \\
V02 & The objective and every research question are approved by the review team, and every object named in a question text has an extraction field owned by that question. & \protocolfield{Pending} \\
V03 & The operational definitions resolve the principal scope ambiguities, and the attribute, deviation, and measure levels are separated and anchored. & \protocolfield{Pending} \\
V04 & The eligibility criteria have positive and negative examples, one inclusion and at most one exclusion criterion per filter, and no filter operating on the review's dependent variable. & \protocolfield{Pending} \\
V05 & Information sources, source retention decisions, and the manual venue sets are validated and documented. & \protocolfield{Pending} \\
V06 & The baseline string retrieves the complete seed set, or every failure has a documented complementary route, and a source contribution matrix exists. & \protocolfield{Pending} \\
V07 & The search string for every source is peer reviewed and preserved verbatim. & \protocolfield{Pending} \\
V08 & Reviewer roles, independence, conflict resolution, and availability are confirmed, and every role concentration is declared. & \protocolfield{Pending} \\
V09 & Both selection calibrations, one for title and abstract and one for full text, are dated and have met the predefined agreement target, and full-text agreement is reported separately from title and abstract agreement. & \protocolfield{Pending} \\
V10 & The mapping, generation-context, and quality-evidence extraction fields are piloted and revised over at least ten studies. & \protocolfield{Pending} \\
V11 & The appraisal instrument is piloted across heterogeneous study designs. & \protocolfield{Pending} \\
V12 & The taxonomy codebook is piloted across more than one UML diagram type, and the unitization rule is exercised by two coders on the same study. & \protocolfield{Pending} \\
V13 & Data management, backup, licensing, and replication package procedures are operational. & \protocolfield{Pending} \\
V14 & The protocol is checked against SEGRESS, PRISMA-P, PRISMA 2020, and PRISMA-S \cite{Kitchenham2023SEGRESS,Shamseer2015PRISMAP,Page2021PRISMA,Rethlefsen2021PRISMAS}. & \protocolfield{Pending} \\
V15 & Manual venue incremental yield and every conditional venue activation decision are documented. & \protocolfield{Pending} \\
V16 & The manual-search provenance chain is operationally tested on at least one core venue-year unit. & \protocolfield{Pending} \\
\bottomrule
\end{longtable}
